% =====================================================================
%  Capitolo 4 — L'atomo di Koster, il production dictionary e la
%  ricerca in pre-produzione (problema e soluzione)
%  Fonte: trascrizione della lezione GD4 (Game Design, primo anno)
%  Pacchetti richiesti nel preambolo del libro:
%    \usepackage[italian]{babel}, \usepackage[T1]{fontenc},
%    \usepackage[utf8]{inputenc} (solo pdfLaTeX), \usepackage{booktabs}
%  Nota sulla fonte: la trascrizione automatica storpia diversi nomi
%  propri (titoli, studi, autori). Dove il riferimento è ricostruibile
%  con certezza lo si riporta nella forma corretta; dove non lo è, lo si
%  segnala in nota invece di indovinare.
% =====================================================================

\chapter[Problema e soluzione]{Problema e soluzione: l'atomo di Koster, il
production dictionary e la ricerca in pre-produzione}
\label{cap:problema-soluzione}

Il capitolo precedente ha aperto il modulo dedicato alla documentazione e al
pitching con il ciclo di produzione, l'iterazione e le regole della buona
documentazione. Questo quarto capitolo lo prosegue e compie un passo
ulteriore: dalla forma dei documenti si passa a \emph{che cosa} si fa prima
ancora di scriverli, cioè alla \textbf{ricerca}. Si torna, in altre parole,
alla ``preistoria'' della produzione, la \emph{pre-produzione}, e a tutto ciò
che, secondo buon uso e costume, andrebbe fatto in quella fase.

La lezione è tenuta ancora da Mattia Podini, in aula a Milano con una parte
della classe collegata da Roma. Il percorso del capitolo ne ricalca la
struttura:
\begin{enumerate}
  \item il concetto di Raph Koster rimasto in sospeso dalla lezione
        precedente: l'\textbf{atomo di gioco}, con il suo ciclo
        problema--azione--feedback--modello mentale
        (sezione~\ref{sec:atomo-koster});
  \item un ripasso rapido della prima lezione di documentazione
        (sezione~\ref{sec:recap-documentazione});
  \item il primo tipo di documento specifico: il \textbf{production
        dictionary} (sezione~\ref{sec:production-dictionary}), seguito da due
        principi generali sulla documentazione, tra cui il fatto che
        \textbf{non esistono template} (sezione~\ref{sec:due-principi});
  \item la nuova lezione, \textbf{problema e soluzione}: il design come
        miglior compromesso possibile, il ruolo dei dati, il \emph{design
        thinking} e la mappa delle attività di ricerca in pre-produzione
        (sezioni~\ref{sec:design-compromesso}, \ref{sec:dati}
        e~\ref{sec:mappa-ricerca});
  \item gli strumenti di quella mappa, uno per uno: competitor diretti e
        indiretti, analisi di utenti, mercato e gameplay, \emph{client pitch
        direction analysis}, stakeholder, problema contro necessità e analisi
        dei competitor (sezioni~\ref{sec:competitor-diretti}--\ref{sec:analisi-competitor});
  \item il task della settimana, il primo production dictionary
        (sezione~\ref{sec:task-gd4}).
\end{enumerate}

La lezione sui problemi e sulle soluzioni non si conclude in questa giornata:
il docente lo anticipa (``non credo riusciremo a finirlo'') e infatti la
chiusura è rimandata alla lezione successiva. Il capitolo riporta quindi la
prima parte, fino all'analisi dei competitor.

% ---------------------------------------------------------------------
\section{L'atomo di gioco di Raph Koster}
\label{sec:atomo-koster}

\subsection{Il punto rimasto in sospeso}

La lezione si apre, come d'abitudine, chiedendo se ci siano domande o dubbi
sulla lezione precedente. Nessuno ne ha; ma, come gli studenti stessi
ricordano al docente, alla fine della volta scorsa era rimasta in sospeso
``una cosa di cui parlare'' (sezione~\ref{sec:task-gd3}). Il docente non
aveva detto quale: lo dice ora.

Si tratta di un'aggiunta alla lezione su che cos'è un gioco
(capitolo~\ref{cap:che-cos-e-un-gioco}), utile anche perché si collega a un
paio di discorsi extra nati con alcuni studenti in fondo all'aula. In quella
lezione il direttore aveva insistito in particolare sul \textbf{mindset
giocoso}, sulla differenza tra \emph{game} e \emph{play}: l'approcciare un
sistema con l'atteggiamento del \emph{play}, su cui si fonda tutto ciò che si
è detto finora e il motivo stesso per cui esiste la disciplina. Noi esseri
umani prendiamo un sistema, vi applichiamo il concetto di \emph{play} e ci
giochiamo sopra. L'esempio era quello di puntare sul mercato azionario; ma il
principio è lo stesso di chi, uscendo dall'aula, vede due piccioni che si
azzuffano e ci scommette sopra (``esempio poco etico, insomma'').
Si veda la distinzione di Koster tra artefatti ludici e sistemi ludici
(sezione~\ref{sec:zuppa}).

Da questa stessa riflessione di Raph Koster deriva un altro concetto: quello
che in lezione viene chiamato \textbf{ludema} o \textbf{game
atom}\footnote{Nella trascrizione il termine compare storpiato
(``l'Udema''). Koster parla di \emph{ludemes} (le unità elementari del
gioco) e di \emph{game atoms}: il nome, come precisa il docente, è poco
importante.}, l'\emph{atomo videoludico}. Il nome importa poco; importa il
concetto, che ritornerà più avanti nel corso ``in diverse salse e in diverse
forme''. Per questo il docente lo presenta adesso: ``ve lo portate a casa''.

\subsection{Gli elementi dell'atomo}

L'atomo di Koster è un piccolo diagramma a ciclo. Lo si percorre punto per
punto, soffermandosi su quelli più rilevanti.

\begin{description}
  \item[Problema] Si parte dal \textbf{problema}, inteso in senso lato (ed è,
        tra l'altro, anche l'argomento della lezione: il termine assumerà più
        avanti un altro significato, sezione~\ref{sec:problema-necessita}).
        Il problema è la \textbf{situazione da affrontare di momento in
        momento} all'interno di un gioco, o di qualsiasi sistema ludico o
        \emph{artefatto ludico}. Il problema sta dalla parte del
        \textbf{sistema}.
  \item[Preparazione] Il problema porta alla \textbf{preparazione} da parte
        del giocatore, che sta dall'altra parte: c'è un boss da battere, come
        mi preparo? Come affronto questo combattimento? C'è una porta da
        aprire, che cosa faccio? ``Gli esempi si sprecano.''
  \item[Core mechanic] Segue una serie di elementi già visti, con
        nomenclature diverse. La \textbf{core mechanic} (meccanica
        \emph{core}) è simile al \emph{verbo di gioco}; ma, come detto la
        lezione precedente, il verbo di gioco nel senso gergale moderno è
        \emph{generico} (sezione~\ref{sec:qa-gd2}), mentre la core mechanic è
        \textbf{specifica}. È ciò che costituisce il gameplay a basso
        livello, l'azione che si ripete tante volte, quella che il giocatore
        compie \textbf{molto più spesso di tutte le altre}. In \emph{Bubble
        Bobble} si sparano le bolle; in \emph{Mario} la core mechanic è
        saltare e muoversi, molto vicina al verbo perché molto semplice.
  \item[Spazio] La core mechanic è sempre \textbf{calata in un contesto}, che
        Koster chiama \emph{spazio} (si potrebbe espandere il discorso, ma per
        ora basta questo). A seconda di dove viene applicata, la stessa
        meccanica cambia effetto. In \emph{Space Invaders} sparare è la core
        mechanic: se tra la propria navicella e l'alieno non c'è niente,
        l'alieno muore; se c'è una copertura, il colpo ne ``mangia'' un pezzo.
  \item[Feedback] Qui si arriva alla parte che interessa di più. Quando si
        compie un'azione, il sistema deve rispondere al giocatore in un dato
        modo: deve dargli il \textbf{feedback}. Ricade qui \textbf{qualsiasi
        cosa comunichi l'effetto desiderato}, o comunque che l'azione è stata
        compiuta. Nel framework di Schell il feedback sta a metà tra
        meccaniche ed estetiche (sezione~\ref{sec:schell}): lo sparo, il
        nemico che esplode, l'effetto sonoro (SFX) che fa \emph{sentire} che
        il nemico è morto, un eventuale punteggio che sale. Sono tutti gli
        esiti, tutti gli effetti di quell'azione. Ciò che il giocatore riceve,
        poi, lo interpreta o meno.
  \item[Successo e fallimento] A seconda della versione del diagramma, a
        questo punto c'è o non c'è una \textbf{biforcazione}. Per introdurre
        il feedback conviene non averla; ma è chiaro che se si compie
        l'azione desiderata, se si ha successo, succede qualcosa, e se si
        compie l'azione ``sbagliata'' succede qualcos'altro. Si apre così il
        \textbf{reame del successo} e il \textbf{reame del fallimento}.
  \item[Ricompensa] Nel reame del successo entra la \textbf{ricompensa}: il
        premio per aver fatto ciò che il designer vorrebbe, quindi qualcosa di
        positivo. In una \emph{boss battle} di \emph{Zelda}: ``hai battuto il
        boss, sei un grande, ti diamo l'arco''. Se si sbaglia succede
        qualcos'altro: il game over.
  \item[Modello mentale] Infine, e qui sta il punto più importante, il
        feedback \textbf{aggiorna il modello mentale del giocatore}.
\end{description}

Il ciclo si può complicare a piacere. In \emph{Dark Souls} ci sono i
\emph{mimic}, i forzieri che in realtà sono mostri: bisogna stare attenti alle
casse. Se si è bravi e si riconosce che non è un mimic, la si apre e si prende
la ricompensa; se si riconosce che è un mimic, la si colpisce, la si sconfigge
e si prende comunque la ricompensa; se non la si riconosce, si viene mangiati.
Reame del fallimento, reame del successo.

\begin{table}[htbp]
  \centering
  \caption{L'atomo di gioco di Koster, come presentato in lezione.}
  \label{tab:atomo-koster}
  \begin{tabular}{@{}p{0.05\textwidth}p{0.22\textwidth}p{0.15\textwidth}p{0.46\textwidth}@{}}
    \toprule
    & \textbf{Elemento} & \textbf{Dove sta} & \textbf{Esempio} \\
    \midrule
    1 & Problema & sistema & un boss da battere, una porta da aprire \\
    2 & Preparazione & giocatore & come affronto questo combattimento? \\
    3 & Core mechanic & giocatore & saltare (\emph{Mario}), sparare (\emph{Space Invaders}) \\
    4 & Spazio & sistema & la copertura che assorbe il colpo \\
    5 & Feedback & sistema & esplosione, SFX, punteggio che sale \\
    6 & Successo / fallimento & sistema & ricompensa (l'arco) o game over \\
    7 & Modello mentale & giocatore & ``i barili rossi esplodono'' \\
    \bottomrule
  \end{tabular}
\end{table}

\subsection{Il modello mentale del giocatore}

Il modello mentale è la parte più importante del framework, anche perché il
resto, osserva il docente, è in parte ripreso da un altro designer, di cui si
vedrà brevemente un altro framework più avanti (e che si ritroverà,
probabilmente, al terzo anno).\footnote{Il nome del designer non è
ricostruibile con certezza dalla trascrizione. L'idea di un ciclo
azione--feedback che aggiorna il modello del giocatore è al centro anche degli
\emph{skill atoms} di Daniel Cook, a cui l'atomo di Koster è spesso
accostato.} Questo elemento sta \textbf{specificamente nel reame del
giocatore}: esiste solo nella sua testa.

Capire di aver fatto qualcosa di positivo o di negativo dipende, tra
virgolette, solo da ``quanto si è svegli''; o meglio, non necessariamente da
quanto si è svegli, ma dal \textbf{bagaglio} che ci si porta dietro. A tutti è
capitato: la prima volta che si prende in mano un controller si è neofiti del
mondo ludico. Il docente ha cominciato, suo malgrado, con \emph{Dino Crisis},
ed era ovviamente ``massacrato ogni due per tre'', tra l'altro da molto
piccolo. Ma con il tempo quel bagaglio cresce e lo si porta con sé.

\paragraph{Le convenzioni codificate.} Nel bagaglio rientrano le convenzioni:
``il barile rosso esplode''. Nella realtà non c'è alcun legame effettivo tra
il colore e l'esplosione (non tutti i barili di benzina sono rossi, né tutti
i barili esplosivi lo sono): è una scelta ormai \textbf{codificata nel mondo
dei videogiochi}, come, banalmente, il forziere. Il giocatore ha imparato
qualcosa del gioco che lo migliora come giocatore e gli permette di
\textbf{approcciare di nuovo il sistema} con informazioni aggiuntive, si
spera da giocatore migliore per quel tipo di gioco.

\paragraph{Il ciclo nella testa del giocatore.} Se si compie un'azione
positiva, il cervello non solo viene nutrito con i feedback giusti, ma
\emph{capisce} di aver fatto giusto: grande scarica di dopamina, e si va
avanti. Se si sbaglia, stesso discorso: forse in quel burrone non ci si
doveva buttare. È ciò che si fa spesso con il \emph{danno da caduta}: esiste
o non esiste? Scopriamolo: ci si tuffa da cinquecento metri e, se si muore,
probabilmente non lo si rifà più. Lo stesso vale per il fuoco amico.

Ecco il legame con la lezione del direttore: la mentalità giocosa, il modo di
approcciare un gioco, è legata al luogo in cui risiede lo stato di
\emph{play}, cioè \textbf{la mente del giocatore}. Lo si espanderà in futuro;
per ora basta sapere che quel bagaglio viene aggiornato continuamente.

\paragraph{Che cosa si può assumere sul giocatore.} A seconda del
\emph{target audience}, si può assumere che il giocatore arrivi nel gioco
con delle competenze, oppure no; in ogni caso \textbf{qualche informazione
bisogna dargliela}. Chi approccia \emph{Dark Souls} da neofita avrà
difficoltà a riconoscere una serie di cose, come i mimic, e potrebbe non
aspettarseli; ma anche un giocatore che non si definirebbe veterano, una
volta mangiato, starà più attento.

\subsection{A che cosa serve nel design}

Una domanda possibile: come aiuta, nel design, tutto questo? In pratica,
\textbf{poco}, se non nel ragionare in un certo reame: quello del problema e
della preparazione, cioè \emph{come presentare il problema} e che cosa il
giocatore deve fare \emph{prima} di affrontarlo. È un discorso complesso, che
per ora resta lì (``un altro strumento misterioso che lasciamo da parte'').

\paragraph{Il playtest.} Ciò che di solito si fa, come accennato la volta
precedente, è \textbf{far testare il proprio gioco ad altri}. Nella propria
testa si ha un'idea di come si sono costruiti il livello, i sistemi e il
resto; ma se si dà il livello in mano a una persona che non ne sa nulla,
esperta o meno che sia, nel 90\% dei casi si scontrerà con cose a cui non si
era mai pensato. La messa in scena e la preparazione verranno trattate in
futuro.

\paragraph{Il feedback non è mai abbastanza.} Il feedback, invece, è
\textbf{importantissimo} e ricade nel reame dell'iterazione visto la volta
precedente (sezione~\ref{sec:principi-iterazione}). Non ci sono mai
abbastanza feedback: si può sempre migliorare il modo in cui si premia o si
punisce il giocatore quando ha fatto qualcosa di giusto o di sbagliato,
anche con modalità che si appoggiano alla psicologia. Il feedback va
\textbf{sempre} curato, spesso manca, e ``non è mai abbastanza''. Si fa il
sistema, si osserva come reagisce il giocatore, si ``pompa'' il feedback o lo
si rende più chiaro (un altro effetto visivo, un SFX migliore) e si ricomincia
il giro: sempre il processo iterativo. Spesso, \textbf{se qualcosa non viene
capito, il problema è lì}.

\paragraph{Approfondimento: la lente del feedback.} Schell dedica al tema
una lente specifica, la \emph{Lens of Feedback} (lente n.~63 nella terza
edizione di \emph{The Art of Game Design}). Il feedback, scrive, è molte cose
insieme: giudizio, ricompensa, istruzione, incoraggiamento e sfida. Per
verificare che il ciclo di feedback crei l'esperienza desiderata, suggerisce
di chiedersi \emph{in ogni momento} del gioco: che cosa i giocatori
\emph{hanno bisogno} di sapere ora? Che cosa \emph{vogliono} sapere? Che cosa
si vuole che \emph{provino}, e quale feedback può creare quella sensazione?
Qual è il loro obiettivo in questo momento? Sono le stesse domande che, nel
ciclo di Koster, collegano il feedback al modello mentale.

\paragraph{Che cosa portarsi a casa.} Più che uno strumento operativo, l'atomo
serve per la \textbf{nomenclatura} e per fissare il concetto di \emph{play}
come stato mentale. Ed è comunque utile avere una parola in più: la
\textbf{core mechanic}, quasi sinonimo di verbo di gioco. L'atomo tornerà
utile più avanti per spiegare perché alcuni design dominanti funzionano così
bene (sezione~\ref{sec:analisi-competitor}).

% ---------------------------------------------------------------------
\section{Un ripasso della documentazione}
\label{sec:recap-documentazione}

Prima di passare oltre, il docente conclude ciò che era stato visto sulla
documentazione la volta precedente, con un ripasso veloce ``perché questa
roba è importante''.

\begin{itemize}
  \item La documentazione è la \textbf{prima \emph{duty}} del designer, il suo
        primo lavoro; a seconda dell'azienda può essere più o meno
        totalizzante (sezione~\ref{sec:perche-documentare}). Serve a
        \textbf{trasmettere la visione}, anche in maniera specifica, con
        documenti adatti (lo si vedrà in futuro); a \textbf{fare ricerca} e a
        documentarla in maniera chiara; a \textbf{supportare il team}; a
        \textbf{controllare la produzione}. Sono i quattro reami principali.
  \item Ci sono le \textbf{golden rules} da applicare
        (sezione~\ref{sec:golden-rules}): ``ricordatevele, stampatevele,
        tatuatevele, fate quello che volete. Ricordatevele, però.''
  \item Ci sono i \textbf{tipi di documenti}. Dal prossimo task questo aspetto
        conterà; per il primo task il docente sarà più permissivo e leggerà
        tutti i documenti, a meno di errori davvero grossolani.
  \item Si usano \textbf{solo documenti della suite di Google} per qualsiasi
        task da consegnare, salvo indicazioni diverse, con le linee guida su
        quale strumento usare in quale momento: \emph{the right tool for the
        right job} (sezione~\ref{sec:scegliere-documento}).
\end{itemize}

\paragraph{La Bibbia di design.} Il discorso della \emph{Bibbia di design}
(detta anche \emph{compendium}) era stato saltato velocemente; se ne
riparlerà quando si tratteranno i tipi di documenti. ``Purtroppo questa piaga
è vera'': esiste ancora oggi, per campi specifici. Di solito questi documenti
sono \textbf{fuori dalla produzione operativa}: non servono a fare il gioco,
ma a questioni, tra virgolette, \emph{politiche}, come rispondere a
stakeholder esterni. Non servono mai a fare effettivamente il gioco.

\paragraph{Documentare qualsiasi cosa.} Con la documentazione si può
comunicare qualsiasi cosa, anche concetti complessi: le \emph{Blueprint} di
Unreal hanno una documentazione che risponde a qualsiasi domanda su un dato
nodo o su una data funzione. In teoria \textbf{tutti i dipartimenti
dovrebbero documentare}; in pratica non sempre lo si fa, e a volte lo si fa
quando la persona non c'è più. Per i programmatori, ad esempio, ``la loro
carne e le loro ossa sono la documentazione'': a meno che un \emph{engine
programmer} non crei un tool che passerà a un altro programmatore che non
incontrerà mai più in vita sua; in quel caso la documentazione la scrive,
altrimenti ``la vedo dura''.

\paragraph{Gli strumenti non sono solo questi.} Quelli indicati nella lezione
precedente sono solo ``briciole'' (sezione~\ref{sec:scegliere-documento}): si
può usare davvero qualsiasi strumento, anche per creare immagini (un
generatore di meme, ad esempio, ammesso che esista ancora).

\paragraph{La dimostrazione pratica.} Il docente assume che i task siano in
corso applicando le regole e le golden rules viste. La dimostrazione
promessa, in cui si prenderà un documento e lo si formatterà da capo, non si
farà in questa settimana, perché la lezione di venerdì sarà da remoto e il
docente preferisce farla ``in carne e ossa'': si farà \textbf{venerdì 6}, la
prima lezione in presenza dopo il 30, compatibilmente con il tempo a
disposizione. Nessuno ha altri dubbi sulla prima lezione di documentazione
(``i dubbi sono sublimati''), e si può passare all'unica cosa rimasta da fare
su quella lezione.

% ---------------------------------------------------------------------
\section{Il production dictionary}
\label{sec:production-dictionary}

\subsection{Un documento comune, con molti nomi}

Nelle lezioni su documentazione e pitching si vedranno diversi tipi di
documenti, ciascuno con il suo scopo. Il primo è il \textbf{production
dictionary}, scelto anche perché è molto comune: probabilmente non è la
prima volta che gli studenti vedono qualcosa di simile. È un
\textbf{dizionario a tutti gli effetti} ed è \textbf{un documento di per sé};
a volte può essere ibridato con altri documenti, ma è davvero raro.

Ha soprattutto molti \textbf{corrispettivi}: documenti che sono la stessa
cosa ``in salsa diversa''. I \emph{narrative designer}, ad esempio (non è
detto che lo si veda nelle lezioni di narrativa, viste le poche ore), lo
chiamano \textbf{lexicon}: vi inseriscono le nomenclature, cioè le singole
voci del dizionario, divise tra nomi comuni e termini che sono neologismi o
idiomi. Il dettaglio interessa poco; il punto da sottolineare è che
\textbf{alcuni di questi documenti hanno un corrispettivo in altri
dipartimenti, o sono comuni a più dipartimenti}. Non è la prima volta che li
si vede: probabilmente è la prima volta che li si trova applicati in una
determinata maniera.

\subsection{A che cosa serve}

Come gran parte della documentazione, il production dictionary serve a
\textbf{spiegare}: il suo obiettivo è \textbf{chiarire nomi specifici}, cioè
la lista di informazioni e nomenclature che il lettore deve conoscere per
comprendere il gioco.

\paragraph{Dove si mette.} Il motivo per cui può essere ibridato è che, a
volte, lo si trova \emph{in testa} (o in coda) ad altri documenti: prima di
leggere il documento, ecco la lista dei termini da sapere. Il docente precisa
di riportarlo perché gli è stato detto, ma di non averlo mai visto in tutti i
suoi anni di lavoro. Il problema è che \textbf{occupa spazio}. Oggi Google
Docs ha le \emph{schede} (\emph{tab}), che permettono soluzioni molto più
eleganti di una prima pagina che è ``un mattone di testo'' sulle informazioni
da leggere. Non è piacevole, soprattutto alla luce delle golden rules: il
documento deve essere esteticamente appagante e deve far venire al lettore la
voglia di leggerlo (sezione~\ref{sec:golden-rules}).

\paragraph{Il libro di cucina.} Si immagini un libro di cucina che, invece di
mostrare subito la grande torta della nonna con le more, con l'immagine e la
ricetta, mostri prima un testo fitto che spiega che cos'è la spatola, che
cos'è la teglia, che cos'è una forchetta, che cos'è un bicchiere. È chiaro
che è meno accattivante per il lettore (``non sono un grande cuoco, questa è
un'altra mia peculiarità''). Di solito, del resto, i termini davvero unici non
sono nemmeno così tanti.

Così, nella stragrande maggioranza dei casi, il production dictionary è un
\textbf{documento a sé}, che funziona come un dizionario e raccoglie tutti i
termini necessari per comprendere il gioco.

\subsection{Che cosa contiene}

Generalmente contiene \textbf{termini unici}: le razze nuove, i tipi di
qualcosa, e simili. Alcuni esempi dalle slide:
\begin{description}
  \item[Risorse con nomi propri] se nel gioco il mana non si chiama
        \emph{mana} ma in un altro modo, va detto. Alla fine dei conti il mana
        funziona come una specifica \emph{currency} che si spende per compiere
        azioni: ``il mana è stamina blu''. Ma serve una definizione, e se
        quella cosa non è mana, va detto che cos'è.
  \item[Termini ormai condivisi] il \emph{mulligan}, oggi condiviso, molti anni
        fa non lo era.\footnote{Nei giochi di carte, il \emph{mulligan} è la
        possibilità di rimescolare o sostituire, del tutto o in parte, la mano
        iniziale.}
  \item[Categorie del proprio sistema] può essere utile definire che cos'è
        una \emph{quest}, se ce ne sono di diversi tipi: se il sistema di quest
        distingue principali, secondarie, terziarie, o prevede delle
        \emph{tier} specifiche, conviene definirle (che cosa significa
        ``primaria'', con un esempio, e così via).
  \item[Stati dell'IA] nelle slide ci sono anche due definizioni di stati
        dell'intelligenza artificiale (IA videoludica, non generativa: se ne
        parlerà nella lezione apposita).
\end{description}

\paragraph{Gli stati dell'IA, un caso esemplare.} A prescindere dal
``cervello'' usato per l'IA (quasi sempre si parlerà di \emph{behavior tree}
o di \emph{state tree}, lo strumento più recente di Unreal Engine 5), cioè da
quel framework applicativo che è lo standard dell'industria, \textbf{i nomi
degli stati in cui si trova l'IA cambiano da azienda ad azienda}. Lo stato
\emph{idle}, o di calma, o qualsiasi nome gli si voglia dare, spesso viene
rinominato. In IO Interactive, ad esempio, uno degli stati dei nemici si
chiama \emph{caccia}; altrove lo si chiama \emph{chase}, altrove ancora in un
altro modo. L'importante è \textbf{far capire a tutti di che cosa si sta
parlando}: è, alla fine dei conti, l'obiettivo della documentazione.

\subsection{Esercizio in aula: una carta di Hearthstone}

Il docente mette alla prova la classe con una carta di \emph{Hearthstone}, un
\emph{Golem cristallino} (che non ha mai giocato, perché è uscito ``da questo
turbinio'' prima che venisse pubblicata): se si dovesse fare il production
dictionary guardando solo questa carta, che cosa bisognerebbe definire?

\begin{itemize}
  \item \textbf{I numeri sulla carta}: costo in mana, attacco, vita. Sono
        termini più o meno condivisi; spiegarli è a discrezione del designer
        o del \emph{lead}, a seconda di \textbf{quanto sono condivisi}.
        Attacco, vita e mana sono generalmente comuni, il che non significa
        che non sia utile spiegarli. Ma conta anche quanto diventa
        ``mastodontico'' il dizionario e quanto tempo di produzione si ha a
        disposizione: come al solito ci sono delle \textbf{priorità}.
  \item \textbf{Le parole chiave}, che a volte non sono parole ma
        \textbf{simboli}.
  \item \textbf{Il tipo della carta}: \emph{Elementale} è una tipologia (una
        sorta di tipo Elettro dei Pokémon).
  \item \textbf{Le keyword dei giochi di carte}, una applicazione specifica,
        vanno spiegate: \emph{Provocazione} (\emph{Taunt}: va attaccato per
        forza prima di poter colpire l'eroe avversario), \emph{Rinascita}
        (\emph{Reborn}: quando muore ritorna in vita), \emph{Rantolo di
        morte} (\emph{Deathrattle}: quando muore produce un certo effetto).
        Sono parole specifiche di \emph{Hearthstone}, non usate altrove.
  \item \textbf{La rarità}: si può andare sempre più in profondità e definire
        anche la rarità. La carta è comune; ci sono i soliti tier (comune,
        rara, epica, leggendaria). \emph{Leggendaria} è già più utile da
        spiegare in \emph{Hearthstone}, perché porta con sé un'informazione
        che il giocatore deve avere: una leggendaria è limitata a
        \textbf{una copia per mazzo}, mentre delle altre se ne possono avere
        più copie.\footnote{In lezione il docente esita tra due e tre copie.
        Nelle regole standard di \emph{Hearthstone} il limite è di due copie
        per le carte non leggendarie e di una per le leggendarie.}
\end{itemize}

\paragraph{La regola: avere priorità.} Bisogna sempre avere una
\textbf{priorità sui termini da spiegare}: quelli più fraintendibili o
necessari per la comprensione vanno spiegati per primi. In questo caso,
\textbf{tutte le keyword}.

Con questo la lezione precedente è conclusa, e il docente assegna il task per
la settimana (sezione~\ref{sec:task-gd4}).

% ---------------------------------------------------------------------
\section{Due principi sulla documentazione}
\label{sec:due-principi}

Una breve parentesi: il docente chiede se gli studenti abbiano recuperato o
rivisto delle lezioni sull'\emph{hub} del corso, e se i paragrafi delle
trascrizioni siano sensati (``non so come la macchina legga queste cose'').
Poi, ``prima che me lo dimentichi di nuovo, con la mia memoria da pesce
rosso'', aggiunge due cose importanti sulla documentazione.

\subsection{La documentazione non sostituisce il gioco}

La documentazione è importante e fondamentale ma, come detto nella lezione
precedente, \textbf{non sostituisce il gioco}. Deve aiutare a creare e
facilitare il lavoro: soprattutto nella pre-produzione, poi
nell'implementazione, nella comunicazione e così via.

\paragraph{Non tutti i documenti sono fondamentali.} Il production dictionary
\emph{potrebbe} essere fatto, ma non è necessario in ogni produzione: come
sempre \textbf{dipende dalle risorse} e dal tempo a disposizione. In una game
jam, con dieci giorni di sviluppo o lavorando in due, non serve; e vale per il
primo documento visto come per molti altri.

\paragraph{I documenti come framework.} Molti di questi documenti, tra
l'altro, possono essere usati \textbf{come framework}: se ne prende un pezzo e
lo si unisce a un'altra parte. La somma di queste cose è il ``principe'' della
documentazione, il \textbf{Game Design Document} (GDD), che il docente spera
di vedere il più avanti possibile (``ma per voi, non per me''). Basta mettere
\emph{DD} alla fine (GDD, LDD\dots) e quello è \emph{il} documento. Dentro ci
finiscono molte cose: il versionamento, potenzialmente il production
dictionary, ``un macello di cose'', a seconda di chi lo scrive e di dove lo si
scrive, cioè in quale azienda. Come al solito, le aziende abilitano o non
abilitano certe cose. Come per i framework, bisogna capire \textbf{perché}
questi documenti si fanno e perché hanno quel contenuto: a volte è un dogma,
più spesso no.

\subsection{Non esistono i template}

La seconda cosa, importantissima, è anche il motivo per cui in questa prima
lezione non c'è un esempio di production dictionary (mentre in futuro ci
saranno esempi di documenti specifici): \textbf{in questo lavoro non esistono
i template}.

Vale soprattutto per i GDD. Cercando corsi, approfondimenti o video, si
trovano persone che condividono il proprio template di Game Design Document.
Va bene, può essere un'informazione utile, e il docente stesso ha trovato
utile vederne alcuni. Ma il consiglio è: \textbf{non applicateli mai uno a
uno}. Anche quando si sarà trovato il proprio metodo di documentazione, e si
faranno GDD in maniera estremamente efficace, non bisogna \textbf{rifugiarsi
dentro template sicuri}, perché non si avrà quasi mai la certezza di poterli
applicare al 100\%: un po' perché i tempi non saranno sempre gli stessi, un
po' perché l'azienda potrebbe non permetterlo (il GDD di alcune aziende ha
quei punti specifici, e quelli sono; e vale anche per altri documenti).

I template aiutano a velocizzare il lavoro, o ad avere in testa che cosa è
importante scrivere in un GDD o in un altro documento. Ma bisogna avere
l'\textbf{elasticità mentale}, e poi la pratica, per cambiarli al volo quando
serve. Se un template viene fornito, vale, ma ha sempre una \textbf{durata
temporale} o, se non temporale, è \textbf{circoscritto all'azienda} in cui si
lavora. Ciò non significa che sia inutile: se ne può prendere il meglio e
ibridarlo con altro, e questa è una grande evoluzione. Ma bisogna
ricordarselo.

\paragraph{Approfondimento: il mito del Game Design Document.} Schell dedica
a questo punto un paragrafo di \emph{The Art of Game Design} (``The Myth of
the Game Design Document''). Molti aspiranti designer, scrive, immaginano che
basti trovare ``il formato giusto'' di GDD per diventare professionisti; e lo
smentisce scrivendolo a caratteri cubitali: \emph{the magic template does not
exist}, il template magico non esiste, non è mai esistito e non esisterà mai.
Riporta anche le parole del designer Jason VandenBerghe: i GDD sono
``letteralmente superati nel momento in cui li si scrive'', perché esprimono
\emph{teorie} su che cosa renderà il gioco buono, non piani o progetti
esecutivi. Ogni gioco è diverso e ha esigenze diverse, di memoria e di
comunicazione, che vanno scoperte di volta in volta: un documento cresce
organicamente intorno al design, non viene da un modello.

% ---------------------------------------------------------------------
\section{Problema e soluzione: il design come compromesso}
\label{sec:design-compromesso}

\subsection{La preistoria della produzione}

Si apre la nuova lezione, \textbf{problema e soluzione}, che parte dalla
``preistoria della produzione'': la \textbf{pre-produzione}. Si tratta tutto
ciò che dovrebbe essere fatto in quella fase secondo buon uso e costume. Non
sempre viene fatto, come si è detto parlando del ciclo iterativo
(sezione~\ref{sec:ciclo-produzione}); ma se venisse fatto, i risultati
sarebbero ``sicuramente molto migliori di quanto non siano nel 50\% dei
casi''.

\subsection{Il miglior compromesso possibile}

Come detto nella primissima lezione (capitolo~\ref{cap:introduzione-game-design}),
\textbf{il design è il miglior compromesso possibile}: fare buon game design
richiede di essere bravi nel compromesso, con le risorse, con la dimensione
del team, con il tempo e così via. Le slide lo riassumono in quattro coppie
di fattori (tabella~\ref{tab:compromesso}).

\begin{description}
  \item[Team e know-how] Quante persone si hanno a disposizione, qual è il
        proprio ``potere di fuoco'', e che cosa si sa fare. Un designer
        esperto di giochi horror verrà tendenzialmente assunto per fare
        quello. Nel mondo \emph{mobile}, guardando le offerte di lavoro, la
        prima cosa richiesta è aver già lavorato nel mobile o nelle
        \emph{live ops} in generale: non significa che chi non lo ha fatto sia
        automaticamente escluso ``da questo magico mondo'', ma si dà priorità
        a chi ha già quel know-how. La \emph{creazione} del team non è
        necessariamente compito del designer; la sua \emph{gestione} sì, è
        possibile.
  \item[Visione e scienza] È la parte \textbf{strettamente legata al
        design}: riguarda \textbf{che cosa il gioco deve essere}. È il
        ``fantomatico'' termine di \emph{visione}, già incontrato rispondendo
        a Lorenzo sui pilastri (sezione~\ref{sec:principi-iterazione}), di cui
        si parlerà meglio con il \textbf{documento di visione}: la
        cristallizzazione dell'idea. Il docente di solito parla di ``messa a
        terra'', ma l'idea non è ancora del tutto messa a terra: ``un piede
        nella fossa e uno no''. Senza parlare ancora di documenti, la visione
        è una \textbf{bussola} che indirizza la produzione: dove si deve
        andare, che cosa si deve raggiungere con il prodotto e perché lo si
        deve fare. E deve essere già stata \textbf{confutata} il più
        possibile tramite il metodo scientifico.
  \item[Mercato e target audience] Il prodotto è indirizzato, creato e
        ragionato sempre in funzione del target audience; il target audience
        è immerso nel mercato, e anche il prodotto, una volta rilasciato, lo
        sarà per forza di cose. Ci si aspetta che l'obiettivo sia vendere, o
        almeno fare \textbf{break even}, cioè guadagnare almeno quanto si è
        speso. Le aspettative possono essere (si spera) molto più idilliache:
        un'azienda non deve solo andare in pari, deve poter ``campare'' altri
        anni. Per farlo bisogna inserire il prodotto in un \textbf{mercato che
        possa assorbirlo}, che abbia abbastanza spazio perché il gioco venga
        visto, giocato e, generalmente, comprato.
  \item[Produzione e soldi] I soldi sono ``la benzina'' dell'industria
        videoludica, almeno ad alto livello (``brutto vederla così, ma è
        così''). Mantenere un individuo costa una certa cifra; i soldi
        permettono di avere un team, di pagare strumenti e ricerca, sondaggi,
        \textbf{focus group} che giochino il gioco per verificare che in un
        certo stato della produzione abbia senso, e aggiustare il tiro (sempre
        l'iterazione). La \textbf{produzione} determina invece quanto è
        efficace la macchina: lavorare 24 ore al giorno o 8, tre giorni su
        sette o cinque, gestire i task in modo che tutti i dipartimenti
        lavorino contemporaneamente o che qualcuno resti fermo. Se i soldi
        sono la benzina, la produzione dice quanto l'ingranaggio è oliato. È
        il lavoro del \emph{producer}, una specializzazione ``tosta'' con i
        suoi documenti specifici da adottare, cambiare e riapplicare, di cui
        non si parla ora.
\end{description}

\begin{table}[htbp]
  \centering
  \caption{Le quattro coppie del compromesso di design.}
  \label{tab:compromesso}
  \begin{tabular}{@{}p{0.25\textwidth}p{0.67\textwidth}@{}}
    \toprule
    \textbf{Coppia} & \textbf{Domanda di fondo} \\
    \midrule
    Team e know-how & Quante persone abbiamo e che cosa sappiamo fare? \\
    Visione e scienza & Che cosa deve essere il gioco, e lo abbiamo
                        verificato con i dati? \\
    Mercato e target audience & A chi lo diamo, e c'è spazio per lui nel
                                mercato? \\
    Produzione e soldi & Quanta benzina abbiamo, e quanto è oliata la
                         macchina? \\
    \bottomrule
  \end{tabular}
\end{table}

\subsection{Cinquanta per cento scienza, cinquanta per cento arte}

Il docente riprende una citazione del direttore, ``molto succinta, molto
\emph{on point}'' (sezione~\ref{sec:inquiry}):

\begin{quote}
\centering
\textbf{Il game design è per il 50\% scienza e per il 50\% arte.}
\end{quote}

La parte \textbf{scientifica}, quella di questa lezione (e non solo), è il
processo di \emph{smontare l'idea}: fare ricerca, trovare i dati, anzi partire
con l'obiettivo di \textbf{avere dei limiti} e di navigare entro quei limiti.
I limiti potenziano la creatività e permettono di evitare strade
problematiche fin dalla partenza. Si vuole fare un MMO, ma si è in tre
persone? Forse non conviene. Si vuole fare un \emph{hero shooter} con gli
\emph{stickman}? Forse sì, forse no: bisogna trovare i dati, non lo si può
dire così. Con il processo scientifico si ottengono dati, li si lima, si
cerca di arrivare a una serie di possibilità, di strade percorribili.

Una volta posto questo filtro scientifico, però, \textbf{dove andare lo
decide il designer}, la sensibilità del singolo: è il 50\% di arte. Può
emergere in momenti diversi della produzione, anche prima, nella scelta
dell'idea; ma l'idea in sé vale poco, perché le idee possono averle tutti.
L'obiettivo è filtrarle, trovare la migliore possibile, applicarla, trovare
l'implementazione giusta. Trovare soluzioni creative fa parte del lavoro, ma
non ne è il tutto: ne è il 50\%.

% ---------------------------------------------------------------------
\section{Senza dati si è solo un'altra persona con un'opinione}
\label{sec:dati}

\begin{quote}
\centering
\textbf{``Senza dati, sei solo un'altra persona con un'opinione.''}\\
\emph{(W.~Edwards Deming)}
\end{quote}

È ciò che serve nel 50\% di scienza del lavoro. Non si può confutare il
design di qualcun altro, e nemmeno iniziare a discuterne, \textbf{senza dati}.
I dati possono essere quelli che il designer ha cercato nella parte di
ricerca del documento, ma vale anche il discorso fatto sul target audience
(sezione~\ref{sec:principi-iterazione}): si può discutere di qualsiasi cosa,
ma la scelta migliore tra la mia e la tua non dipende da quanto è forte l'uno
o l'altro; dipende da \textbf{chi avrà in mano il gioco}, da chi dovrebbe
giocarlo.

\paragraph{Gioco ``bello'', gioco ``brutto''.} È il discorso fatto con
Federica (sezione~\ref{sec:qa-gd2}): non serve dire che un gioco è bello o brutto (parole bandite,
sezione~\ref{sec:vocabolario}); serve dire che tipo di estetiche ha. Sono dati
\emph{qualitativi}, ma sono dati, e sono verificabili. Lo stesso vale per
``all'utenza piacerebbe'' o ``non piacerebbe'': si fa ricerca, si chiede
all'utenza, al potenziale gruppo di persone. Se bisogna fare un gioco che
venda su Reddit, si va su Reddit a chiedere; se servono i più grandi fanatici
dei \emph{soulslike}, nel peggiore dei casi si fermano le persone per strada:
``sei un gran fan dei soulslike?''. Prima o poi si trova qualcuno, gli si fa
una serie di domande, si raccolgono i dati, li si filtra. Servono per iniziare
a parlare di questioni concrete, o per andare a discuterne.

Deming, economista statunitense, ha fatto molto altro: ha avuto grande
successo in Giappone, in particolare sui \emph{processi produttivi}, temi più
legati alla sfera della produzione che a quella del design; magari se ne
parlerà più avanti.

\subsection{Quanti dati, e quando}

Come si ottengono i dati? ``Guarda un po''': tramite la \textbf{ricerca}, la
prima parte di questa lezione. Si possono raccogliere dati
\textbf{quantitativi} e \textbf{qualitativi}, in qualsiasi maniera:
l'importante è che \textbf{siano dati}.

Chiedere a una o due persone del target audience se ha senso fare una cosa
o un'altra non basta: due persone sono troppo poche per un bacino statistico,
ne servono \textbf{almeno cento}. Non sempre si avrà il tempo di farlo; ma
nella pre-produzione la parte di ricerca va fatta \textbf{attentamente},
proprio per questo. Si costruiscono fondamenta solide, si raccoglie una
``marea'' di dati che magari non servono tutti subito ma a cui si può tornare
in futuro, e soprattutto, facendolo ora, non lo si deve fare nel vivo della
produzione, quando il tempo scarseggerà davvero.

La ricerca si fa in diverse maniere e forme: alcune molto aziendali e
condivise, altre meno. La lezione serve proprio a dare \textbf{più strumenti
possibili} per la ricerca e, potenzialmente, per l'ideazione e il
\emph{brainstorming}. Non tutti saranno sempre applicabili: è una lista ``più
che estensiva'', e va valutato di volta in volta che cosa usare.

\subsection{Il design thinking in cinque passi}

Per rendere concreto il discorso, il docente mostra un breve video sul
\textbf{design thinking}, che descrive il processo come un viaggio senza
destinazione certa, ``ma con la fiducia che si arriverà in un posto
fantastico''. L'esempio: un team riceve la sfida di far muovere di più i
bambini americani, per contrastare l'obesità infantile. L'idea di partenza è
dare ai bambini un lettore musicale digitale con un sensore di movimento,
premiandoli in base all'attività. Ma la grande domanda è: i bambini lo
userebbero davvero, e abbastanza a lungo da vederne i benefici?\footnote{Il
video è quello con cui lo studio di design IDEO illustra il proprio approccio
al design thinking; il progetto descritto ha portato a un dispositivo di
monitoraggio dell'attività per ragazzi, collegato a una piattaforma online.}

\begin{enumerate}
  \item \textbf{Imparare dalle persone} (\emph{learn from people}). Il team
        parla con i bambini, passa del tempo nelle loro case e nelle loro
        scuole, ascolta i loro limiti, le abitudini, ciò che li entusiasma e
        ciò che li frustra. La ricerca include bambini ``nella media'', ma
        anche bambini molto attivi e molto sedentari: chi sta agli
        \textbf{estremi} sa dare voce a problemi che chi sta nel mezzo
        percepisce, ma fatica a mettere a fuoco.
  \item \textbf{Trovare degli schemi} (\emph{find patterns}). Le osservazioni
        finiscono su centinaia di post-it, disposti per dare un senso a ciò
        che si è imparato; con un'intuizione informata si cercano gli schemi
        che indicano delle opportunità. Daniel dice: ``mi annoio con i
        videogiochi in solitaria, sono quelli multigiocatore che mi fanno
        tornare''. Meg dice: ``non porto l'iPod quando corro perché voglio
        parlare con le persone''.
  \item \textbf{Definire i principi di design} (\emph{define design
        principles}). Frasi come queste rivelano uno dei principi che
        porteranno al concept: \emph{facilitare l'interazione sociale in ogni
        momento}. Ne emergono molti, e insieme formano i ``pali della porta''
        di un'esperienza che il team è fiducioso possa funzionare con i
        bambini.
  \item \textbf{Rendere tangibile} (\emph{make tangible}). Domande del tipo
        ``come potremmo\dots?'' colmano la distanza tra principi e idee
        specifiche; le migliori vengono trasformate rapidamente in prototipi
        grezzi.
  \item \textbf{Iterare senza sosta} (\emph{iterate relentlessly}).
        Dispositivi fisici di cartone ed esperienze digitali simulate con
        carta e matita permettono di imparare in fretta; a ogni prototipo il
        concept si evolve, viene riportato ai bambini, si ascolta e si impara.
\end{enumerate}

Il risultato non è l'idea iniziale del lettore musicale, ma un piccolo
\textbf{monitor di attività} da agganciare ai vestiti o da mettere in tasca:
un portale verso un mondo online in cui i ragazzi condividono e celebrano con
gli altri e con le famiglie i risultati raggiunti nel mondo reale. In uno
studio clinico, il risultato è stato un \textbf{aumento del 59\%
dell'attività fisica}.

``Grande successo: l'America è salvata dall'obesità'', commenta il docente.
L'obiettivo era far alzare i bambini dalla sedia e ``fargli toccare l'erba''.
È un esempio molto contenuto di cinque punti che possono essere espansi, e
che sono contenuti nella prima parte della lezione: sono \textbf{il cuore} di
ciò che si vedrà, e non vanno dimenticati.

\subsection{Partire da un problema o da una necessità}

Il primo passo è partire da un \textbf{problema} o da una \textbf{necessità}
(la distinzione verrà chiarita nella sezione~\ref{sec:problema-necessita}):
può essere semplicemente qual è l'obiettivo del progetto, del team. A volte su
questo obiettivo si costruisce un'identità intera, cosa più riconoscibile tra
i publisher, ma visibile anche tra gli studi.

\paragraph{Esempio: Milestone.} Quale potrebbe essere la missione di
Milestone? Per dirlo bisogna sapere che cosa fa: \emph{MotoGP}, i giochi su
Valentino Rossi e simili. L'obiettivo è fare \textbf{simulazioni di guida},
giochi a due e quattro ruote simulativi. Magari lo studio non era nato con
quell'obiettivo, ma nel tempo si è solidificato: hanno già chiaro che cosa
devono fare, quale problema risolvere e in quale mercato affacciarsi. Per
un'altra azienda può non valere: si potrebbe volere ``un sacco di soldi'', e
allora l'obiettivo è trovare un'applicazione, un progetto che dia quel
risultato. Come? Facendo ricerca: chiedendo alle persone se c'è un problema o
una necessità.

\subsection{Esercizio in aula: che cosa non va nei soulslike?}

Per mostrare come si raccolgono i dati, il docente intervista i giocatori di
soulslike presenti in aula: che cosa non piace dei soulslike? Le risposte:
\begin{itemize}
  \item una studentessa, gran fan del genere, fatica a trovare difetti; l'unica
        cosa che le viene in mente riguarda la \textbf{console} su cui gioca,
        non il gioco;
  \item Pedro cita le \textbf{risorse limitate}: in alcuni titoli, finiti gli
        oggetti curativi, bisognava tornare a \textbf{farmarli} prima di
        riprovare (il \emph{grinding});
  \item Antonino cita la \textbf{curva di difficoltà} tra l'inizio e il resto
        del gioco: l'inizio è difficilissimo;
  \item un altro studente porta l'esempio del DLC di \emph{Elden Ring}: la
        difficoltà principale dei boss, secondo lui, non nasce dall'averli
        resi più dinamici, ma dall'aver alzato i numeri (vita, velocità,
        attacco, danno), mentre nei capitoli precedenti il combattimento era
        un'alternanza bilanciata tra schivata e attacco, sia del giocatore sia
        del boss.
\end{itemize}

Il docente avverte che il discorso è complesso e lo si sta semplificando.
Se l'obiettivo è fare un soulslike, bisogna capire come ci si è arrivati: un
buco di mercato, giocatori che chiedono qualcosa a gran voce, mille altre
ragioni (``magari siamo FromSoftware e facciamo quello''). E bisogna chiedere
al proprio target audience che cosa vuole.

\paragraph{La curva a campana.} Va anche capito se gli intervistati sono
\textbf{esattamente il tipo di giocatore che si vuole}, e dove si collocano
sulla curva. La distribuzione normale vista nel video è la condizione
naturale delle cose: se si chiede a tutta l'Italia ``ti piace la pizza?'', nel
mezzo della curva a campana ci sono quelli che la mangiano una volta a
settimana, e agli estremi chi non la mangia mai e chi la mangia ogni singolo
giorno. Bisogna capire dove si trovano gli intervistati: grandissimi fan di
\emph{Dark Souls}, con ``un santino di Miyazaki a casa'', oppure no. Al
docente, ad esempio, dei soulslike non piace per nulla la narrativa,
``estremamente criptica, convoluta, complicata''; ma non sono quelli i
problemi emersi in aula.

\paragraph{Registrare tutto.} Dalle risposte emergono già due questioni,
\textbf{complessità/difficoltà} e \textbf{grinding}. Queste interviste servono
a prendere dati: bisogna segnare tutto e, potenzialmente, \textbf{filmare}
(la lezione stessa è registrata). Bisogna osservare le reazioni delle
persone: se fanno fatica a trovare le risposte, se si animano parlando di
qualcosa. Non conta solo ciò che dicono, ma \emph{come} lo mostrano e che
cosa potrebbero pensare. Poi si setacciano i dati per vedere quali tornano
utili.

\paragraph{Trovare archetipi e categorie.} Il modo che conviene più di tutti è
\textbf{trovare archetipi}, categorie in cui far ricadere le risposte. È anche
per questo che è stata data la lista di emozioni (sezione~\ref{sec:vocabolario}),
e il concetto verrà esplorato e rifinito. Tutto ciò di cui si è parlato in
aula (la progressione, il tipo e la curva di progressione, la sfida, quanto
alta o bassa, il grinding) ricade nello stesso reame: la \textbf{challenge}.
Uno studente osserva che in \emph{Elden Ring} il problema degli oggetti
curativi è stato risolto (le fiasche si ricaricano ogni volta che si muore),
segno che probabilmente un processo del genere c'è stato; e lo stesso si può
dire della telecamera e del \emph{level design} dei primi \emph{Dark Souls}.
Se invece qualcuno avesse detto ``in \emph{Dark Souls} mi piace un sacco
parlare con il Cavaliere Cipolla'', o ``vorrei un compagno, vorrei un
cane'', ci starebbe (``chi sono io per dire di no''), ma sarebbe rimasto fuori
da questa potenziale ``molla'' della challenge, che include sfida, grinding e
progressione.

\subsection{Correlare, chiudere la visione, prototipare}

Una volta ottenuta questa ``enorme bacinona di dati'', bisogna cercare la
\textbf{correlazione}: trovare gruppi, cose che collidono, capire \emph{perché}
succedono certe cose. È più evidente, ad esempio, quando un mercato è molto
saturo e le persone non vogliono più giocare a un certo genere: la questione
del recente filone dei \emph{friendslop}, per esempio.\footnote{Termine
gergale per l'ondata di giochi cooperativi economici, pensati per essere
giocati con gli amici in streaming. Nella trascrizione compare come
``friend slot''.} Oppure il motivo per cui una parte dei giocatori dedica ai
videogiochi meno tempo che, poniamo, nel 2022: il Covid ha avuto la sua
influenza. Bisogna trovare questi legami per poter rispondere sensatamente
alle domande che ci si pone dentro il team, e anche fuori.

\paragraph{Chiudere la visione.} Verso la metà o la parte finale della
ricerca bisogna, se non lo si è già fatto, \textbf{iniziare e poi chiudere la
visione}. Ciò che si deve fare con il progetto viene \textbf{blindato}:
questi sono i punti, queste le cose incluse, che per forza di cose ci devono
essere. Altre, per forza di cose, non ci saranno. Se si fa un soulslike, i
giocatori si aspettano un certo tipo di difficoltà e certi sistemi: se si
muore si perdono le risorse e bisogna andarle a riprendere, è lo standard. A
volte è un'aspettativa anche la narrativa, la cosiddetta narrativa
\emph{apofenica} dei soulslike (motivo per cui il docente ``se la sorbisce''
anche in \emph{Hollow Knight}). La visione è ciò che deve condividere tutto il
team: tutti devono essere \emph{onboarded}, devono sapere dove si va e qual è
la direzione del progetto.

\paragraph{Prototipare.} Poi il \textbf{prototyping}: prima lo si fa, meglio
è (lo si è già visto, sezione~\ref{sec:segreti-iterazione}). Si raggiungono il
prima possibile versioni concrete, fisiche o digitali, per testare ciò che si
sta facendo; le si gioca, si raccolgono feedback, si modifica, si itera. In
pre-produzione lo si fa in maniera ridotta rispetto alla produzione, ma già si
mettono le mani nell'engine: è sicuramente ancora pre-produzione.

% ---------------------------------------------------------------------
\section{La mappa della ricerca in pre-produzione}
\label{sec:mappa-ricerca}

Il docente presenta i punti che verranno toccati nella prima parte della
lezione. Come sempre, la lista \textbf{non è esclusiva né limitata} a queste
voci: sono modi di fare ricerca, prendere dati specifici e potenzialmente
correlarli.

\begin{description}
  \item[Client pitch direction analysis] Capire se si ha \textbf{qualcuno
        sopra}, se c'è una situazione ``clientelare'': i limiti non sono solo
        quelli del target audience, ma anche quelli di chi fornisce i soldi.
        Se si fa un progetto per qualcun altro, bisogna capire che cosa vuole
        quella persona; e, come sempre, \textbf{ciò che le persone dicono non è
        sempre ciò che vogliono}: va tradotto, contestualizzato, analizzato.
        È legata alla definizione di stakeholder
        (sezioni~\ref{sec:client-pitch} e~\ref{sec:stakeholder}).
  \item[Stakeholder e il loro peso] Non tutti quelli che stanno sopra pesano
        allo stesso modo. Un investitore che ha messo un milione di dollari
        nel gioco ma ha il grosso dei suoi interessi altrove (``nel campo del
        petrolio'') si interesserà relativamente del progetto: qualsiasi cosa
        si faccia, bene o male andrà in porto (``la semplifico molto, ma
        molto''). Uno stakeholder che si aggiorna spesso, impone deadline
        specifiche ed è sempre ``sulla spalla'' a controllare ha un peso
        diverso. C'è chi non dà soldi e non influisce economicamente, e chi
        dà soldi senza avere competenze definite in ambito videoludico.
        Bisogna capire chi c'è sopra, che cosa vuole e quale peso ha nella
        produzione.
  \item[Problema contro necessità] Capire se c'è un problema o una
        necessità da risolvere nel mercato, e non solo
        (sezione~\ref{sec:problema-necessita}).
  \item[Analisi dei competitor] Fondamentale: ``non si sfugge''. Là fuori ci
        sono giochi che hanno qualcosa in comune con il proprio, anche se il
        proprio non è ancora uscito. Quali sono i propri competitor,
        \textbf{diretti} o \textbf{indiretti}? (Sezioni~\ref{sec:competitor-diretti}
        e~\ref{sec:analisi-competitor}.)
  \item[Analisi degli utenti] Chi dovrebbe giocare il gioco
        (sezione~\ref{sec:utenti-mercato}).
  \item[Analisi di mercato] Dove ci si affaccia, e con quali regole.
  \item[Analisi di gameplay] Correlare i dati e mettersi la maschera del
        target audience, fino alla \textbf{visione}.
\end{description}

% ---------------------------------------------------------------------
\section{Competitor diretti e indiretti}
\label{sec:competitor-diretti}

\subsection{Il competitor diretto}

Per chiarire la distinzione, il docente chiede un genere alla classe: i
soulslike. Mostra \emph{Enotria: The Last Song} e chiede quale potrà mai
essere il competitor diretto di un soulslike (``perché avete paura di
parlare, ragazzi?''). La risposta è ovvia: \textbf{qualsiasi cosa faccia
FromSoftware}. Sono competitor diretti perché condividono la stragrande
maggioranza degli aspetti del gioco: il sistema di progressione, la
difficoltà, potenzialmente la telecamera.

Come osserva Federica, però, a seconda di quanto è innovativo un gioco
\textbf{i competitor diretti non sono mai uno solo}: bisogna individuare quelli
più vicini. Lo scopo del progetto e la sua grandezza sono molto minori
rispetto a una produzione FromSoftware, quindi il competitor più diretto sarà
qualcosa di più indipendente (un ``doppia A'', mettiamola così); ma il gioco è
immerso nel mercato dei soulslike, e per forza di cose \emph{Dark Souls} sarà
uno dei suoi competitor diretti. Artisticamente, però, \emph{Enotria} è
diverso: gli sviluppatori lo definiscono un \textbf{\emph{Summer Souls}}, un
soulslike ``estivo''. Anche questa è una scelta di vendita, come si era detto
dei \emph{roguelite} nella prima lezione.

\subsection{Il competitor indiretto}

Si può trovare anche un \textbf{competitor indiretto}: un gioco che ha in
comune con il proprio \emph{solo} un aspetto, artistico o meccanico, una
specifica meccanica, un sistema o una direzione artistica. Qualcosa di
\textbf{estremamente limitato e circoscritto} che ricorda un altro gioco, con
cui però non si compete direttamente nello stesso mercato.

\paragraph{Pokémon e Palworld.} \emph{Palworld}, o qualsiasi titolo simile,
è un competitor \emph{diretto} di \emph{Pokémon}, perché va a intaccarne il
monopolio. Ha un'ampia distanza su una serie di cose (``lo schiavismo dei
Pokémon, le armi, le mitragliatrici''), ma lo stile artistico è molto simile
a quello dei \emph{companion} e dei \emph{monster collectible}, il genere su
cui Pokémon ha una presenza estrema, un vero monopolio di mercato.
\emph{Clair Obscur: Expedition 33}, invece, non lo è.

\paragraph{Un caso italiano.} Il docente cita anche un \emph{monster
collector} di un team italiano, che riprende da Pokémon e da altri titoli il
sistema di combattimento con la \emph{timeline} e ha una direzione artistica
che sembra animata da uno studio d'animazione giapponese.\footnote{I titoli e
gli studi citati in questo passaggio non sono ricostruibili con certezza dalla
trascrizione; si riporta il ragionamento, non i nomi.} Non possono competere
con Pokémon: hanno un costo in risorse diverso (e bisognerebbe capire quanto
investono Nintendo e Game Freak) e uno scopo diverso, perché il loro è,
teoricamente, un \emph{roguelike}. Pokémon resta un competitor diretto, ma
hanno \textbf{deviato il tiro} in un'altra direzione, per intercettare un
certo tipo di persone ed evitare certi problemi, o rispondere a certe
necessità, previste dalla ricerca. Anche questo è il punto della ricerca:
identificare i competitor diretti o simili aiuta a capire che cosa gli altri
hanno fatto bene o male, se si può fare meglio o peggio (``peggio si evita,
chiaramente'') e soprattutto \textbf{come ci si può differenziare}. E poi i
competitor vanno \textbf{analizzati}: non restano lì nel nulla.

\subsection{Gli esempi della classe}

La classe propone diversi esempi, e la discussione è vivace.
\begin{description}
  \item[Pokémon TCG Pocket e Yu-Gi-Oh!] Un gioco mobile, quindi bisogna
        capire il panorama mobile e prendere i TCG mobile con partite di
        durata ridotta. Indiretti potrebbero essere \emph{Hearthstone} o
        \emph{Gwent}; diretti, probabilmente, \emph{Yu-Gi-Oh! Duel Links} (e,
        in misura minore, \emph{Master Duel}), forse \emph{Marvel Snap}.
  \item[Expedition 33 e la serie Mario \& Luigi] Uno studente li avvicina per
        il controllo che il giocatore ha durante il turno nemico (il
        \emph{parry}, le azioni a tempo): non è un gioco di ruolo a turni
        ``puro''. Il docente concorda: è \textbf{indiretto}. Non lo fa solo
        lui, ma c'è quella specifica cosa in comune. Allo stesso modo il
        sistema di morte di \emph{Dark Souls} viene ripreso in titoli che non
        sono soulslike. Pokémon ed \emph{Expedition 33} hanno meccaniche
        simili, ma giocatori ``completamente opposti''.
  \item[Sekiro e Final Fantasy VII Remake] Un altro studente si chiede se
        siano indiretti per il combattimento reattivo, non strategico. Non
        sono paragoni sbagliati: il punto è che \textbf{è uno spettro}, e
        quanto ci si avvicina o allontana dipende.
  \item[Resident Evil 2 e Silent Hill 2] Probabilmente diretti, forse più in
        passato che oggi: \emph{Resident Evil} si è spostato molto più verso
        l'action che verso il survival, mentre \emph{Silent Hill}, pur non
        avendo mantenuto del tutto la visione iniziale, resta molto meno
        action.
\end{description}

\paragraph{Indiretto non significa ``stesso target, mezzi diversi''.} Uno
studente chiede se indiretto significhi che il target è lo stesso ma con
mezzi diversi. \textbf{Non necessariamente}. C'è un elemento, o più di uno, in
comune con il proprio prodotto, ma è estremamente limitato; a volte è
proprio un mercato diverso. Chi condivide lo stesso sistema ha molto più in
comune di un \emph{Expedition 33}; e due giochi con arte completamente diversa
ma meccaniche simili, o con la stessa arte ma meccaniche diverse, sono casi
ancora diversi. ``Sono tutti esempi giusti'', conclude il docente: sono su uno
spettro, su scale diverse, ed era lui a sbagliare insistendo a cercarne uno
ancora più vicino.

\paragraph{Un caso di confine: Nine Sols e Hollow Knight.} Questi due giochi
stanno davvero \textbf{sulla linea di demarcazione} tra diretto e indiretto.
Possono essere percepiti come molto simili (e andrebbe verificato se i
giocatori di \emph{Hollow Knight} sono andati a comprare \emph{Nine Sols} per
la somiglianza); ma a livello di gameplay e di combattimento sono molto
lontani: \emph{Nine Sols} è molto più basato sul \emph{parry} (il docente ci
fa ``una fatica tremenda''), \emph{Hollow Knight} è più veloce. Sono lontani
anche nella trama: in \emph{Nine Sols} la parte narrativa, con il sistema di
dialoghi e i personaggi che parlano, è molto più spinta ed esplicita. Senza
dati, il docente li considera indiretti, ``ma proprio di un pelo''.

\subsection{Come individuarli: i punti insindacabili}

Come sempre è uno spettro, e lo si verifica \textbf{con i dati}: si fa
ricerca, si prende una lista di più giochi possibile. Un modo giusto per
arrivarci è partire da ciò che per il proprio gioco è
\textbf{insindacabilmente necessario}: le cose ad alto livello che il gioco
deve avere, ad esempio le \emph{game qualities} di Romero e Schreiber
(sezione~\ref{sec:romero-schreiber}). Magari sono già state messe sul tavolo
dagli stakeholder: se si deve fare un monster collectible si è già molto
dentro certi limiti. Se invece si deve fare un generico gioco da tavolo, bisogna
capire l'obiettivo: se è vendere tanto, forse conviene guardare ai
\emph{party game}, ai giochi popolari che hanno venduto oltre una certa soglia,
ai più conosciuti, o a quelli entro una certa fascia di prezzo.

Trovati i punti insindacabili, si limita su di essi la ricerca dei
competitor:
\begin{itemize}
  \item i giochi che ne hanno \textbf{di più} in comune, e sono più vicini, sono
        i \textbf{competitor diretti};
  \item quelli che ne hanno magari \textbf{uno solo}, in maniera estremamente
        attinente, sono gli \textbf{indiretti}.
\end{itemize}
Non è un discorso facile di base, e serve tanta ricerca per far emergere
queste cose; lo si approfondirà.

\subsection{L'opposizione}

Si possono identificare anche i giochi in \textbf{opposizione}: titoli molto
lontani dal proprio, in netta contrapposizione. A che cosa servono? Non sempre
a qualcosa: può non valere per prodotti che puntano a un mercato specifico. Ma
possono servire a \textbf{fare innovazione}, ad esempio unire una direzione
artistica unica a un gameplay mai fatto prima, o a qualcosa che nessuno ha
fatto (o, molto più probabilmente, che nessuno ha fatto con successo). Unire
titoli o meccaniche davvero lontani tra loro può \textbf{stimolare la
creatività}, anche nel brainstorming.

Il docente prova a pensare a un esempio: un action con un'arte molto poco
leggibile, come \emph{Superhot} (``il tempo si muove solo se ti muovi''), che
ha un'arte simile a quella di \emph{The Unfinished Swan}. Non essendo un
artista, nota però alcune differenze: \emph{The Unfinished Swan} è più
morbido nelle linee, più circolare, \emph{Superhot} è molto più spigoloso; ma
nella palette, bianco e nero, ``siamo lì''.

\paragraph{That's Not My Neighbor e Papers, Please.} Uno studente chiede se
\emph{That's Not My Neighbor} e \emph{Papers, Please} siano indiretti. Sì, per
il docente è un buon esempio. Il primo è un gioco abbastanza semplice in cui
si controllano le persone che bussano alla porta prima di farle entrare, con
una deriva horror; \emph{Papers, Please} è radicato in una realtà fittizia,
con la sua dinamica burocratica. Lo studente osserva che sono giochi di una
nicchia strana, in cui è difficile trovare un competitor diretto. È un'altra
domanda da porsi: perché hanno meccaniche simili? Perché è una nicchia di
mercato estremamente stretta, e chi gioca quel tipo di giochi può richiedere
necessità specifiche. Probabilmente \emph{That's Not My Neighbor} è andato in
un'altra direzione proprio per non scontrarsi direttamente.

\paragraph{I competitor nel pitch.} Soprattutto nel \emph{pitching}, i
competitor rendono molto più facile comunicare, anche la commistione di
elementi che non sono meccaniche: ``il nostro gioco è come \emph{Limbo}, ma
l'arte è completamente diversa''. Si chiama in causa un gioco che magari non
c'entra nulla, ma ha un punto artistico o sonoro vicino al proprio. Come
analizzare i punti tra competitor diretti e indiretti lo si vedrà; e quasi
certamente lo tratterà anche Daniele nelle lezioni di narrativa, perché il
marketing è la sua grande specialità.

\paragraph{Analizzare i propri giochi.} Valerio chiede se, potendo disporre di
una vasta varietà di giochi propri (ad esempio, per Pokémon, i titoli
principali e quelli secondari), li si possa usare nell'analisi. Certamente: si
parte avvantaggiati, un po' perché la documentazione c'è già, un po' perché,
nel caso di Game Freak, si ha il monopolio del mercato e si fa quella cosa da
più di vent'anni. Proprio per questo diventa \textbf{ancora più importante
chiedere al target audience che cosa vuole}; e quello di Pokémon non è uno
solo, è molto frammentato. Ciò che il docente chiede a Pokémon non è ciò che
chiedono i ragazzini che lo giocano per la prima volta, e il rapporto è
probabilmente di uno a dieci. È il motivo per cui \emph{Scarlatto e Violetto}
e \emph{Spada e Scudo} hanno una difficoltà e una struttura così lineari ed
esplicite: la frangia di veterani è in numero molto minore rispetto a chi
deve comprare il gioco. Si potrebbe anche discutere del fatto che i giochi
non siano la principale fonte di guadagno di The Pokémon Company, che vuole
vendere soprattutto le carte; ma per vendere peluche e carte serve una nuova
generazione di Pokémon.

% ---------------------------------------------------------------------
\section{Analisi degli utenti, del mercato e del gameplay}
\label{sec:utenti-mercato}

\subsection{Analisi degli utenti}

Come si cercano i giochi simili al proprio, si cerca di capire \textbf{chi
dovrebbe giocare} il proprio gioco. È la ``cartina tornasole'' di tutte le
scelte di produzione: ciò che si fa è legato a questi dati. Se sono sbagliati,
o estremamente vaghi, diventa difficile ragionare sulla soluzione migliore,
perché non si è stretta abbastanza la ``falce'' sulle possibilità: si
potrebbe fare qualsiasi cosa. Bisogna quindi capire:
\begin{itemize}
  \item la \textbf{demografica} degli utenti;
  \item che cosa fanno, che cosa giocano;
  \item che cosa piace loro e che cosa non piace;
  \item che cosa dicono dei titoli già rilasciati e, in futuro, del proprio.
\end{itemize}

Ci si può già chiedere se un'idea funziona: si va dal proprio target
audience ideale e si chiede ``ti piacerebbe se facessimo \emph{Dark Souls}
con i pony di \emph{My Little Pony}?''. Se l'80\% delle persone risponde
``incredibile, non vedevo l'ora'', si sta andando nella direzione giusta; se
risponde ``non so\dots'', allora no.

\paragraph{Il vantaggio dell'opposizione.} È anche il vantaggio di cercare i
giochi in opposizione: unire generi lontani anni luce. Il docente giocherebbe
volentieri \emph{My Little Pony} in versione \emph{Dark Souls}, e fa notare
che sono giochi che esistono. \emph{Dark Souls} è \emph{grimdark}, è
\emph{Berserk}, è l'atmosfera cupa di un mondo completamente distrutto e alla
deriva; se se ne prendono le meccaniche e lo si porta in fondo al mare, in un
mondo da cartone animato, si ottiene \emph{Another Crab's Treasure}, un
soulslike con un granchio come protagonista. E, sempre per il target
audience, ha una \emph{easy mode} in cui il granchio ha una pistola e spara a
tutti i boss: un pubblico chiaramente molto lontano da quello di FromSoftware,
o meglio, non solo quello. \emph{Palworld} funziona esattamente così:
l'innocenza dei Pokémon unita allo schiavismo delle multinazionali e alle
armi.

\subsection{Analisi di mercato}

Dove ci si affaccia? In quale mercato? È \textbf{saturo} o no?
\begin{itemize}
  \item un \textbf{blue ocean} è uno spazio in cui c'è qualcosa che nessuno ha
        mai fatto (chiedendosi perché);
  \item un \textbf{red ocean} è il contrario: un mercato estremamente saturo
        di prodotti, in cui è estremamente difficile emergere. Lo sono, ad
        esempio, i \emph{battle royale}, gli \emph{hero shooter}, i
        \emph{friendslop}, e ce ne sono altri.
\end{itemize}

\paragraph{Approfondimento: oceani e barriere.} Schell riprende la stessa
metafora in \emph{The Art of Game Design}: i mercati inondati di offerta sono
\emph{red ocean}, un mondo ``pesce mangia pesce'' con molto sangue
nell'acqua, in cui la sfida non è più arrivare nello store ma emergere tra
migliaia di concorrenti. Per sopravvivere o si è il pesce più grande, o si
trova un \emph{blue ocean}, un posto con pochi concorrenti: a volte perché il
mercato è nuovo, a volte perché ci sono \textbf{barriere} che tengono fuori
gli altri (tecniche, hardware, di competenza, di vendita e distribuzione). Il
suo esempio: in un mercato saturo come quello delle app iOS, all'epoca, circa
il 90\% dei giochi perdeva soldi; nel mercato delle slot machine, protetto da
regolamenti, hardware e competenze specifiche, circa il 90\% era in attivo.

\paragraph{Le regole del mercato.} A seconda del mercato a cui ci si affaccia,
alcune cose si possono fare e altre no, perché ci sono regole.
\begin{itemize}
  \item \textbf{Cina}: se l'obiettivo è rilasciare in Cina, bisogna
        premunirsi. O si fa un gioco senza teschi, senza ossa e senza
        iconografia legata alla morte o al cristianesimo (come le tombe con
        la croce), oppure si dovrà rifare il lavoro dopo. È successo, ad
        esempio, a \emph{Rainbow Six Siege}, che per il rilascio in Cina ha
        dovuto cambiare numerosi asset.
  \item \textbf{Germania}: non si può avere iconografia nazista; la svastica
        non può comparire in nessuna forma.
  \item \textbf{Belgio}: lo stesso vale per le \emph{loot box}.
\end{itemize}
O ci si pensa prima, decidendo di non rilasciare in quel mercato, o si
produce in modo che la produzione abbia quel margine. Se nel gioco ci sono
ossa o tombe con la croce, e quell'asset si può rendere generico (versione X,
versione Y, cambia solo una \emph{decal} o un nome), adattarlo è semplice; se
invece si è costruito un livello con mille asset raggruppati, diventa
difficile. Poi ci sono casi come \emph{Dead Space}: i nemici sono
\emph{necromorfi}, bestie mutate che si disintegrano con le ossa che escono,
grondanti di sangue. \emph{Dead Space} non si può fare senza necromorfi, e
quindi non si vede in Cina: ma ci si preclude una grossa fetta di mercato.
Questo ragionamento \textbf{va fatto prima}.

\subsection{Analisi di gameplay e visione}

L'analisi di gameplay, generica, si vedrà alla fine di tutto: si
\textbf{correlano i dati} per raggiungere una certa consapevolezza. È come
\textbf{mettersi la maschera del target audience}. Oltre ad avere dati su cui
tornare e su cui discutere, l'obiettivo è mettersi nei panni di chi giocherà
il gioco, capirne le necessità, i desideri, ciò che gli piacerebbe fare.

È più facile se si fa un gioco per cui si è già naturalmente predisposti come
giocatori, ma va fatto comunque. Ed è più difficile, ed è ciò che fa
diventare professionisti, mettersi nei panni di qualcuno molto lontano da ciò
che si gioca. Anche il primo caso ha le sue insidie: il docente, grande amante
dei TCG, quando ci ha lavorato ha sbagliato, mettendo \textbf{il proprio ego
dentro il design}: voleva più complessità, perché è ciò che si aspetterebbe
lui da giocatore, ma non necessariamente ciò che si aspetta il target
audience di quel gioco.

La concretizzazione di tutto questo è la \textbf{visione di gioco} (\emph{game
vision}), molto più generale dei pilastri. Si cristallizza in un documento
che, quando verrà fatto (e quando il docente lo insegnerà), ``quello è e
quello rimane'': \textbf{la visione non può cambiare in corsa}. Ci sono dei
limiti, e quelli sono.

\subsection{Dopo la pausa: game qualities e generi}

Dopo una pausa di sei minuti (``ci rivediamo alle 4''), e qualche battuta con
Alessandro, studente del terzo anno presente in aula per studiare, il docente
risponde a una domanda arrivata durante la pausa sulle \emph{game qualities}
come punti insindacabili. Se le game qualities sono ``prima persona, PvP'',
i giochi sono un milione. Se si aggiunge il genere, \emph{hero shooter}, la
falce si chiude: dentro ci sono \emph{Overwatch}, \emph{Deadlock},
\emph{Valorant} (in prima persona), \emph{Team Fortress}. Identificare altri
elementi, che non fanno necessariamente parte delle game qualities ma sono
più specifici, rende più semplice spiegare il gioco e anche
\textbf{pitcharlo}. Ci può stare anche la direzione artistica, se ad alto
livello; ma bisogna considerare che le game qualities sono comunque
\textbf{molto macroscopiche}.

% ---------------------------------------------------------------------
\section{Client pitch direction analysis}
\label{sec:client-pitch}

Si torna alla mappa e la si percorre punto per punto, partendo dalla
\emph{client pitch direction analysis}: i modi per cercare, fare domande,
chiedere.

\subsection{Fate domande}

All'inizio si fanno \textbf{tante domande}. Ma vale sempre, sottolinea il
docente: se si hanno dubbi, si fanno domande; se c'è qualcosa, si parla, si
chiede. In aula e fuori, non bisogna avere paura di fare domande, perché
altrimenti il rischio è \textbf{darsi da soli una risposta sbagliata}. Vale
anche al lavoro: se non si capisce qualcosa, si chiede; se viene fatta una
domanda senza tutte le informazioni necessarie, si chiede, ``per favore,
parlate'' (si veda anche la sezione~\ref{sec:soft-skill}).

\subsection{Le cinque W}

I metodi sono estremamente semplici. Il primo sono le classiche domande
inglesi, le \textbf{cinque W} (\emph{where}, \emph{when}, \emph{who},
\emph{why}, \emph{what}), a cui si aggiunge \emph{how}.\footnote{In aula la
lista viene ricostruita a memoria con l'aiuto degli studenti; la formula
classica è \emph{5W + 1H}.} Danno informazioni specifiche su una serie di
elementi molto diversi, come lo spazio e il tempo. Non sono un framework: sono
domande, la base per spiegare qualsiasi tipo di storia. Se si deve spiegare il
concept del proprio gioco, lo si può fare con queste domande in una frase, o
al massimo in cinque righe, e diventa molto chiaro:
\begin{itemize}
  \item \emph{dove} è ambientato il gioco?
  \item \emph{quando} è ambientato?
  \item \emph{chi} è il protagonista?
  \item \emph{perché} fa ciò che sta facendo?
  \item \emph{che cosa} deve fare?
  \item \emph{come} lo fa?
\end{itemize}
Le risposte iniziano a separare le cose. \emph{Metroid} e \emph{Hollow
Knight}, pur avendo lo stesso sistema di progressione, hanno stili di
combattimento diversi: il \emph{what} e lo \emph{how} cambiano notevolmente
(e, tra i due, anche il \emph{who}).

\subsection{I nove perché e la root cause analysis}

L'altro metodo sono i \textbf{nove perché}: qui le domande sono solo
\emph{why}. C'è un ``mantra'' secondo cui si può raggiungere la radice di
qualsiasi problema rispondendo nove volte alla domanda ``perché?''. C'è un
problema: ci si chiede perché; l'interlocutore dà una risposta; si chiede di
nuovo perché, e così via. Aiuta a \textbf{eliminare possibili
fraintendimenti}. È il metodo dei bambini molto piccoli, che continuano a
chiedere ``perché?''; filosoficamente si potrebbe dire che serve a raggiungere
l'essenza delle cose (il docente non ci crede, ma la possibilità c'è).
Elencare le risposte a questi perché, sempre più in profondità, è la
\textbf{root cause analysis}, l'analisi delle cause
radice.\footnote{Nella letteratura di management la tecnica è più nota come
\emph{5 Whys} (``cinque perché''), nata nei processi produttivi di Toyota; la
versione a nove domande, \emph{Nine Whys}, è una delle \emph{Liberating
Structures} citate più avanti (sezione~\ref{sec:problema-necessita}). Il
principio è lo stesso: ripetere la domanda finché non emerge la causa di
fondo.}

\subsection{Osservare, intervistare, filmare}

\begin{itemize}
  \item \textbf{Chiedere}: se c'è qualsiasi dubbio, se qualcosa è formulato in
        maniera incomprensibile, si chiede.
  \item \textbf{Osservare e intervistare} l'utente e il cliente, con
        \textbf{domande strutturate}. Quella fatta in aula ai fan dei
        soulslike (``che cosa ti piace, che cosa non ti piace?'') è una
        domanda possibile, ma estremamente vaga. Con \emph{poll} e
        \emph{survey} (se ne parlerà) si possono limitare le risposte, se si ha
        già un'idea di dove ci si vuole muovere, o migliorare la domanda
        andando sempre più nello specifico.
  \item \textbf{Filmare} le persone, se possibile, per non perdere dati:
        mentre giocano, mentre le si intervista. Non bisogna perdersi niente:
        anche i responsi emotivi, i sorrisetti, i dubbi, i silenzi sono
        importanti. È un materiale enorme e non è detto che serva tutto; ma
        se serve, è lì.
  \item \textbf{Giocare e interagire}: tutto questo dovrebbe portare
        all'\emph{insight}; ci si arriverà, e insieme alla correlazione porterà
        ad altro.
\end{itemize}

\subsection{Priorità e visione del cliente}

\paragraph{Ordinare le priorità.} Il cliente avrà priorità maggiori e
minori. A volte consegna la lista di ciò che vuole, ma \textbf{ordinarla spetta
al designer}: alcune cose saranno più pertinenti, altre più negoziabili.

\paragraph{La visione e la linea editoriale del cliente.} Qual è la visione
del cliente? Come si diceva di Milestone, il cliente può avere una
\emph{missione}, una verticalità costruita negli anni: fa solo giochi, o
videogiochi, o solo un certo tipo di esperienze. Se si cercano finanziamenti
per un gioco con un mago gatto che spara vodka infuocata, forse a
\textbf{Devolver Digital} può interessare: il suo catalogo è pieno di giochi
di questo tipo, come \emph{Enter the Gungeon}, dove tutto è proiettili, o
\emph{My Friend Pedro}, dove una banana parlante guida il protagonista a
suon di sparatorie. Le \textbf{linee editoriali} sono importanti: il cliente
potrebbe avere già una visione specifica da abbracciare, se si lavora sotto di
lui.

\subsection{Must have e nice to have}

Questo è l'esempio che si riporterà nella visione: le cose da fare per forza.
\begin{description}
  \item[Must have] Ciò che è \textbf{indiscutibile}, senza compromessi: le
        meccaniche, i sistemi, le cose da implementare per forza.
  \item[Nice to have] Ciò che è buono che ci sia, ma non è detto che ci sarà.
        Generalmente sono \textbf{aggiunte alla produzione} a cui,
        realisticamente, ci si può dedicare se si hanno più tempo e risorse,
        dopo aver risolto i must have. Se mancassero, l'anima del gioco, il suo
        \emph{core}, resterebbe intatta: il gioco resta solido, comunica lo
        stesso messaggio, resta coerente con sé stesso. ``Se ci fossero
        sarebbe meglio; se non ci sono, vanno in pace, ce ne faremo una
        ragione.''
\end{description}

\paragraph{Dove finiscono i nice to have.} Uno studente cita un gioco uscito
senza un sistema di mappa, aggiunto due mesi dopo il
lancio.\footnote{Il titolo non è ricostruibile con certezza dalla
trascrizione.} Esatto: \textbf{DLC e patch} sono gli esempi perfetti per
inserire i nice to have o per rispondere a esigenze dell'utenza a cui prima
non si era pensato o che non si riuscivano a risolvere. Nel DLC di \emph{Elden
Ring}, ad esempio, l'utenza faticava a battere Radahn, e il sistema è stato
modificato in risposta (si può discutere se sia stato giusto, se il feedback
sia stato recepito e se la reazione sia stata corretta, ma è successo). Uno
studente aggiunge che, secondo lui, anche la seconda fase del boss finale del
gioco base era pensata per essere affrontata a cavallo, cosa che non è
successa; il docente non lo sapeva, ma è ``un ottimo esempio''. Sono sistemi
molto specifici che restano nel reame del ``possiamo farlo se'': se si fa
tutto il resto, se si ha tempo, se si hanno le risorse. Vanno identificati e
messi nel documento di visione; se ne riparlerà.

\subsection{Esempi concreti e situazioni estreme}

\paragraph{Usare esempi concreti.} Mentre si parla con l'utente o con il
cliente, si usano esempi di giochi simili per capire se la linea desiderata è
quella: giochi di successo e giochi che hanno fallito. Più avanti si cercherà
anche di capire \emph{perché} hanno fallito: se il proprio competitor diretto
è un gioco che ha fallito, bisogna assolutamente rispondere a quella domanda.
L'assunto di base, chiaramente, è che il proprio gioco non fallirà.

\paragraph{Pensare a situazioni estreme.} Se nessuno giocasse al nostro gioco,
falliremmo? Va capito. Non è necessariamente compito del designer, ma se si
dovesse fare un piano di produzione e di rilascio, ci sono diverse fasce:
\begin{itemize}
  \item quante copie servono per arrivare al \textbf{break even}, in pari tra
        spese e guadagni?
  \item in quante copie si traduce un \textbf{insuccesso}? E un risultato
        \textbf{disastroso} (probabilmente zero clienti)?
  \item quanto vale un \textbf{successo}? E un \textbf{successo folle}, quello
        per cui ``andiamo tutti alle Maldive per tutta la vita'' (o dovunque si
        voglia andare)?
\end{itemize}
E se si avesse un milione di utenti? Se il gioco è online, con server
dedicati, e l'affluenza è esagerata, che cosa succede? È capitato più volte:
con \emph{Overwatch} all'inizio, e di recente con un altro gioco \emph{live
service}. Bisogna pensarci prima, per \textbf{limitare i danni}: prevedere
tutto non è possibile, ``è inumano''. E se la Terra esplodesse?
Provocatoriamente: forse si giocherebbe su Marte (``io sicuramente non ci
giocherei, perché non ci arriveremo''). Anche questo è un esempio di root
cause analysis.

% ---------------------------------------------------------------------
\section{Gli stakeholder}
\label{sec:stakeholder}

\subsection{Chi sono}

Uno \textbf{stakeholder} è qualsiasi persona, organizzazione, gruppo sociale
o società in senso lato che ha uno \emph{stake}, una posta, nel progetto: ha
qualcosa da dire, ``ha le mani in pasta''. Ha dei \textbf{diritti} sul
progetto, ma si porta dietro anche dei \textbf{doveri}. Gli stakeholder vanno
\textbf{identificati}: bisogna capire, dentro e fuori dal team, chi c'è e
quanta importanza ha.

Anche gli utenti rientrano in questa categoria; ma di solito, parlando di
stakeholder, è più rilevante \textbf{chi sta sopra} che chi sta sotto. Gli
utenti sono i ``principi'', ma qui interessa chi può modificare la
traiettoria del progetto al di là di ciò che si potrebbe prevedere, chi ha (o
avrebbe) qualcosa da dire durante una deadline, una milestone e così via. Si
definisce, in sostanza, \textbf{tutto ciò che è fuori dal team operativo}.

Si torna alla piramide già vista: i \emph{director} (creative director, game
director, art director, sound director\dots), poi i \emph{lead}, poi le
seniority (senior, mid, junior, come si vuole). Accanto c'è la produzione, con
il \emph{producer}, già un livello a sé, che potrebbe essere identificato come
stakeholder; e da un'altra parte, completamente, c'è il \textbf{team di
marketing}.

\paragraph{Il marketing e il CEO.} ``Una cosa molto brutta che succederà a voi
e succederà anche a me in futuro'': il team di marketing può bussare alla
porta e dire ``questa cosa non è abbastanza \emph{screenshottabile}''. Il
designer può rispondere che a lui serve la meccanica, non che finisca su
Instagram; ma il marketing può comunque venire a fare richieste. Oppure c'è il
\textbf{CEO}, che con la produzione non ha nulla a che vedere: ``ieri ho visto
su History Channel che Giulio Cesare aveva una spada gigantesca: mettiamo una
spada gigantesca nel gioco''. Non dovrebbe essere il designer a rispondere di
no: la richiesta dovrebbe essere filtrata prima, ma a volte non succede. Il CEO
è uno stakeholder \emph{definitivo}, come si vedrà.

\subsection{Classificarli}

Ci sono diversi modi di identificare gli stakeholder e di dar loro un volto
\textbf{archetipale}, cioè farli ricadere in gruppi specifici. Aiuta a
ragionare e a rispondere in maniera diversa, ma codificata, a seconda del
gruppo: al marketing in un modo A, al CEO in un modo B, e così via.

\paragraph{La matrice potere--interesse.} Un primo modo è collocarli su un
piano cartesiano con due assi, \textbf{potere} e \textbf{interesse}
(tabella~\ref{tab:potere-interesse}).

\begin{table}[htbp]
  \centering
  \caption{La matrice potere--interesse degli stakeholder.}
  \label{tab:potere-interesse}
  \begin{tabular}{@{}p{0.2\textwidth}p{0.36\textwidth}p{0.36\textwidth}@{}}
    \toprule
    & \textbf{Interesse alto} & \textbf{Interesse basso} \\
    \midrule
    \textbf{Potere alto} & \emph{Promotori} (es.\ il CEO): tenerli buoni,
      aggiornarli e rassicurarli costantemente &
      \emph{Latenti} (es.\ un investitore interessato solo ai ricavi):
      tenerli soddisfatti \\
    \addlinespace
    \textbf{Potere basso} & \emph{Difensori} (es.\ la critica
      videoludica): tenerli informati &
      \emph{Apatici}: monitorarli, perché possono cambiare posizione \\
    \bottomrule
  \end{tabular}
\end{table}

\begin{description}
  \item[Promotori] molto potere e molto interesse: il CEO, ad esempio. Vanno
        tenuti buoni sempre, aggiornati costantemente ed eventualmente
        rassicurati, costantemente.
  \item[Latenti] molto potere ma poco interesse: vanno \textbf{tenuti
        soddisfatti}. Se un'azienda ha investito nel gioco ma non è
        interessata a come esce, bensì a quanti soldi fa, le si danno
        informazioni sulla produzione: se si è in pari con le deadline, se i
        costi sono nel reame del previsto, se i prospetti di vendita sono
        allineati. Se si è cambiata una meccanica, se si è aggiunto o tolto un
        boss, se era un drago gigante o una viverna sottomarina, al latente
        interessa poco; e non conviene nemmeno dirglielo, perché potrebbe
        creare più problemi che altro. Se c'è un problema di costi, si parla di
        quello. Il design è meno rilevante.
  \item[Apatici] poco potere e poco interesse: vanno \textbf{monitorati}.
        Possono cambiare il loro valore sulla scala: in questo momento della
        produzione restano lì, ma in futuro potrebbero avere più potere o più
        interesse.
  \item[Difensori] poco potere e molto interesse: vanno \textbf{tenuti
        informati}. Possono essere, ad esempio, la critica videoludica (e
        bisognerebbe capire quanto potere abbia nel mondo moderno:
        relativamente poco, ``o almeno mi piace pensarlo''). Non hanno potere
        sulla modifica del progetto, ma vanno tenuti informati.
\end{description}

\paragraph{Il modello della salienza.} Un altro modo usa un diagramma a tre
insiemi\footnote{Si tratta del \emph{modello della salienza degli
stakeholder} (\emph{stakeholder salience model}) di Mitchell, Agle e Wood; nella
trascrizione il nome è storpiato.}, le cui categorie sono \textbf{potere},
\textbf{legittimità} e \textbf{urgenza}. Uno stakeholder può avere molto
potere sul progetto ma poca legittimità (ciò che dice è poco attinente), e
molta urgenza (deve essere aggiornato di frequente). Un producer, ad esempio,
ha generalmente sempre molta urgenza e, non necessariamente ma a seconda dei
casi, meno legittimità nelle sue richieste. Lo stakeholder
\textbf{definitivo} le ha tutte e tre; gli altri ne uniscono due o ne hanno
una sola. Non si entra troppo nel dettaglio per questioni di tempo, ma lo
strumento è nelle slide.

% ---------------------------------------------------------------------
\section{Problema contro necessità}
\label{sec:problema-necessita}

\subsection{Due definizioni}

Nel gergo comune tutto è un problema: se c'è qualcosa da affrontare, un
intoppo, è un problema. Ma si può dare una definizione più sottile, che
distingue il problema dalla necessità.

\begin{description}
  \item[Problema] appare in una \textbf{struttura già esistente}. È qualcosa
        di facile da risolvere, ed è necessario risolverlo in maniera
        efficiente: prima lo si risolve, meglio è. Un buco nel muro; nascondere
        ``la polvere sotto il tappeto''; qualcosa che si fa velocemente. Un
        \textbf{bug} è generalmente un problema, se non nasce da una struttura
        disintegrata del sistema. Di solito si assume (il docente lo assume)
        che il programmatore pensi l'implementazione in prospettiva futura:
        rende esplicite le variabili per i designer, o ne nasconde alcune
        perché il designer non distrugga tutto mettendo uno zero dove non
        deve. Se in un sistema del genere capita un bug, è un problema. Nel
        reame della vendita: se i fan dei soulslike chiedono qualcosa per
        sistemare la difficoltà o per avere una curva di progressione diversa,
        è un problema, perché sta in una struttura già esistente (il soulslike
        stesso): vanno toccati alcuni punti che non modificano l'identità del
        titolo.
  \item[Necessità] è qualcosa di \textbf{nuovo}, che può portare a grossi
        danni a cascata. Se là fuori qualcuno chiede a gran voce un monster
        collectible in cui si schiavizzano i mostri, li si ``massacra di
        botte'' e li si arma con un AK-47, quella è una necessità: non appare
        in una struttura esistente, bisogna \textbf{costruirci un gioco
        attorno}.
\end{description}

\paragraph{Non sono mutualmente esclusivi.} Uno studente riprende il boss
finale di \emph{Elden Ring} pensato per essere combattuto a cavallo, dove il
cavallo non si poteva usare nelle prime patch: problema o necessità? Si può stare \textbf{a metà tra le due}. Probabilmente era più vicino alla necessità, ma è
stato risolto come un problema, per i soliti limiti di tempo, soldi e
produzione. Nulla vieta di farlo, se il tempo è poco.

\paragraph{Un caso vissuto: il quest log.} Al docente è capitato di avere
sistemi complessi che, tramite l'iterazione, mostravano delle mancanze: un
sistema di quest, o il fatto che in un certo momento della progressione
dovesse accadere un evento specifico. Serviva una \textbf{struttura nuova},
proprio per il discorso sulla programmazione fatto prima: qualcosa che
accogliesse un sistema di quest complesso, un \emph{quest log} fatto come si
deve, sia nella presentazione all'utente sia ``sotto il cofano'', con
un'implementazione sensata. Non c'era tempo, e si faceva un
\textbf{``accrocchio'' in Blueprint}, qualcosa che funzionasse e che, lato
utente, desse lo stesso risultato: la quest veniva risolta, la porta si
apriva, l'oggetto veniva dato. Ma non era la forma corretta: lì serviva un
sistema complesso, era una necessità, e per necessità (appunto) la si è
trattata come un problema. Non è l'ideale, ma capita. Succede anche quando si
fa un titolo innovativo o si vuole testare un sistema nuovo: è una necessità
perché richiede qualcosa che non è mai stato fatto, o un lavoro complesso che
si porta dietro una forma complessa.

\paragraph{Forma e funzione.} In modo semplicistico: nella necessità
\textbf{la forma è maggiore della funzione} (quanto maggiore, va capito); nel
problema la funzione ``ha la meglio''. È la differenza tra mettere un cerotto
e fare i punti di sutura, o asportare un rene e richiudere la ferita. Sistemi
complessi, potenzialmente innovativi, richiedono emozioni complesse da
evocare e pensare, meccaniche e sistemi da soddisfare.

\paragraph{Approfondimento: la lente del problem statement.} Schell parte
dallo stesso punto: lo scopo del design è risolvere problemi, e prima di
generare idee bisogna sapere \emph{perché} lo si fa. Un buon \emph{problem
statement} dichiara insieme l'obiettivo e i vincoli (``come posso fare un
gioco per browser che piaccia davvero agli adolescenti?''), e va calibrato:
troppo largo porta a soluzioni che non centrano l'obiettivo, troppo stretto
taglia fuori soluzioni brillanti. I vantaggi sono tre: uno spazio creativo più
ampio, un criterio chiaro per misurare le idee (quanto bene risolvono il
problema?) e una comunicazione migliore nel team. La \emph{Lens of the Problem
Statement} (lente n.~14) chiede: quale problema sto davvero cercando di
risolvere? Ho fatto supposizioni che non c'entrano con il vero scopo? Un gioco
è davvero la soluzione migliore? Come saprò che il problema è risolto? In apertura
del paragrafo Schell riporta una massima anonima: ``non innamorarti della tua
soluzione, innamorati del tuo problema''.

\subsection{Strumenti per metterlo a terra}

I modi per mettere a terra questa analisi sono parecchi.

\paragraph{Il diagramma a lisca di pesce.} Il \textbf{diagramma di
Ishikawa} (a lisca di pesce) aiuta a identificare le cause secondo
dipartimenti, materiali, macchine e così via: per ogni area principale si
individuano cause primarie e secondarie, disposte ``a lisca''. È un modo tra
tanti, lasciato a disposizione.

\paragraph{Le Liberating Structures.} Molti di questi strumenti si trovano
nelle \textbf{Liberating Structures}, una raccolta di metodi che si usano sia
nella fase di ricerca sia, più avanti, nel brainstorming. Esistono anche in
forma di carte, con un valore simile a quello delle carte di design di Schell
(il \emph{Deck of Lenses}, sezione~\ref{sec:lenti}), per risolvere problemi.
Tra queste ci sono i \textbf{nove perché} e, ad esempio, la \emph{Troika
Consulting}, una consulenza a turno in piccolo gruppo. Altri strumenti si
vedranno in futuro, sempre per il brainstorming.

\paragraph{Il framework Cynefin.} Un altro modo, se non si sbaglia derivato
dallo stesso ambito\footnote{Si tratta del framework \emph{Cynefin}, ideato
da Dave Snowden, che le Liberating Structures includono tra i propri
strumenti.}, invece di distinguere problema e necessità, categorizza
``l'impiccio'' che si ha davanti (per non dire problema) in quattro categorie:
\textbf{semplice}, \textbf{complicato}, \textbf{complesso},
\textbf{caotico} (tabella~\ref{tab:cynefin}).

\begin{description}
  \item[Semplice] La difficoltà sta \textbf{nel compiere l'azione}, mai nella
        lista di azioni da fare. Non saper fare una UI in Unreal può essere un
        problema semplice: si cerca un tutorial, si ha la lista di cose da
        fare, le si fa, risolto. Una ricetta: fare la torta è un problema
        semplice, perché si assume che, seguendo gli step esattamente come
        indicato, il risultato sia \textbf{altamente riproducibile}.
  \item[Complicato] Il grado di difficoltà si alza, perché la lista di
        protocolli non è legata a una sola disciplina, ma a \textbf{più
        discipline}, magari molto lontane tra loro e molto complesse. Far
        decollare un razzo richiede calcoli matematici complessi, simulazioni
        fisiche, materia prima (una lega che resista a quel calore),
        l'addestramento di chi ci sale, il carburante. La percentuale di
        successo è molto più bassa che nel semplice: \textbf{si può
        prevedere}, si può ridurre il rischio, come si fa con qualsiasi lancio,
        ma è più difficile.
  \item[Complesso] La prevedibilità viene meno: il problema diventa
        \textbf{non categorizzabile}. L'esempio è educare un bambino, perché
        ogni individuo è a sé. Chi ha fratelli o sorelle è stato
        probabilmente educato in modo simile, perché il modello educativo dei
        genitori è quello; ma a seconda dei casi i genitori aggiustano il
        tiro sulla risposta dell'``utente'', il figlio. Il docente ha una
        sorella ed è completamente diverso da lei, nonostante siano cresciuti
        in maniera molto simile.
  \item[Caotico] L'unica discriminante la fanno la \textbf{collaborazione}
        tra i membri e l'\textbf{improvvisazione}. Gli esempi sono mosca cieca
        e \emph{Taboo}: un certo numero di parole da non dire, un tempo
        limitato per far indovinare una parola al proprio compagno di squadra.
        Tra amici si usano \emph{inside joke}, un gergo che solo il gruppo
        conosce: far capire qualcosa a qualcuno con cui si ha grande intimità è
        più semplice, perché si condividono esperienze e intuito. Per spiegare
        una parola si cita un meme che entrambi conoscono.
\end{description}

\begin{table}[htbp]
  \centering
  \caption{Le quattro categorie dell'``impiccio'' (framework Cynefin).}
  \label{tab:cynefin}
  \begin{tabular}{@{}p{0.16\textwidth}p{0.42\textwidth}p{0.34\textwidth}@{}}
    \toprule
    \textbf{Categoria} & \textbf{Dove sta la difficoltà} & \textbf{Esempio} \\
    \midrule
    Semplice & nell'eseguire, non nel sapere che cosa fare & una ricetta, un
      tutorial di UI in Unreal \\
    Complicato & più discipline intrecciate; prevedibile ma rischioso & il lancio
      di un razzo \\
    Complesso & non prevedibile né categorizzabile; si aggiusta il tiro &
      educare un bambino \\
    Caotico & solo collaborazione e improvvisazione & mosca cieca,
      \emph{Taboo} \\
    \bottomrule
  \end{tabular}
\end{table}

Sono modi per \textbf{categorizzare} il problema, la necessità, qualsiasi cosa
ci sia da affrontare: dargli un nome o una categoria può aiutare, per
semplicità.

% ---------------------------------------------------------------------
\section{Analisi dei competitor}
\label{sec:analisi-competitor}

\subsection{La SWOT e le domande}

Ci sono diversi modi, più o meno precisi, per analizzare i competitor, e
molti si vedranno nel corso: lo stesso framework dei \emph{core elements}
(sezione~\ref{sec:schell}) serve anche all'analisi, ed è proprio ciò che si
sta facendo nel task; non è l'unico.

\paragraph{La SWOT analysis.} Tra i modi ormai conosciuti c'è la
\textbf{SWOT analysis}: punti di forza (\emph{strengths}), punti di
debolezza (\emph{weaknesses}), opportunità (\emph{opportunities}) e minacce
(\emph{threats}), cioè ciò che minaccia il progetto. La si può disporre a
matrice. Non si usa solo nell'industria videoludica, anzi: è presa dal
management generico.

\paragraph{Le domande sui competitor.} Per i singoli competitor conviene però
porsi delle domande, soprattutto sui \textbf{diretti}, che vanno analizzati
in profondità (``su questo non ci piove''):
\begin{itemize}
  \item che cosa fanno, e come lo fanno?
  \item c'è un \textbf{design dominante}? Perché?
  \item c'è qualche design che si tende a \textbf{evitare}? Perché?
  \item quanto hanno venduto?
  \item quali sono i loro \textbf{unique selling point}?
  \item che cosa piace al loro target audience, e che cosa no? Che cosa
        comunica? Ci sono parole specifiche?
  \item c'è una categoria specifica di \emph{influencer} che adotta o gioca
        quel tipo di giochi? (Ormai è necessario per il marketing.)
  \item ci sono feature di gameplay specifiche che vengono mostrate, o no, nel
        marketing? Come si fa marketing, con quale tono di voce? (Questo
        ``grosso peso'' il docente lo lascia a Daniele.)
  \item sono stati un successo o un fallimento? Perché?
  \item che cosa si può fare realisticamente, sempre entro soldi, team e
        così via?
\end{itemize}

\paragraph{Informarsi senza giocare.} Bisogna guardare interviste, talk di
settore, come hanno fatto. E se non si riesce a giocare un titolo (magari non
è ancora uscito) si guarda il \textbf{gameplay senza commento}, come diceva
il direttore. Nemmeno il docente gioca a tutto, e realisticamente non
potrebbe: ci sono modi laterali per ottenere informazioni su un gioco
guardandolo.

\subsection{Il design dominante}

Il \textbf{design dominante} è ciò che sta dietro al detto ``non reinventate
la ruota'': quella ruota esiste già da qualche parte. È qualcosa che viene
riapplicato $N$ volte, nella stessa identica maniera: perché?

In \emph{Dark Souls}, ad esempio, il sistema di morte e recupero delle anime
è uno standard, e non solo lì: anche fuori dal genere. C'è un motivo: è
\textbf{estremamente efficace}. E qui torna l'atomo di Koster visto
all'inizio (sezione~\ref{sec:atomo-koster}): \emph{Dark Souls} funziona così
bene proprio per quel discorso di \textbf{feedback} e di \textbf{educazione
del modello mentale del giocatore}. Se ne riparlerà quando si parlerà di
psicologia, a fine anno.

\paragraph{Si può brevettare una meccanica?} Uno studente chiede se si possa
brevettare una meccanica. Il grande vantaggio del game design è che
\textbf{le meccaniche non hanno brevetti}, con un asterisco. Sono come
\textbf{blocchi di Lego}: così specifiche nell'applicazione, ma così generali
nelle regole. Non si può brevettare il salto, perché il verbo di gioco è lo
stesso tra \emph{Mario} e \emph{Uncharted}, ma le applicazioni sono
completamente diverse; anche tra \emph{Celeste} e \emph{Mario}, pur essendo
entrambi platform. Si possono quindi prendere meccaniche e sistemi di gioco
uno a uno da giochi anche lontani dal proprio, e va benissimo. A volte lo si fa
per un gioco intero: il primo \emph{Lords of the Fallen} è, di fatto, ``la
\emph{skin} di \emph{Dark Souls}'', con gran parte dei sistemi riprodotti uno
a uno. Funziona; non è particolarmente creativo e potrebbe non essere
l'aspirazione di un designer, ma si fa, per tempistiche e per semplicità.

\paragraph{I casi limite.} Ci sono casi limite di brevetti, ma sono
estremamente rari.
\begin{itemize}
  \item \textbf{Pokémon}: il brevetto usato nella causa contro \emph{Palworld}
        riguarda, se il docente ricorda bene, la creatura che finisce dentro
        un oggetto portatile e la possibilità di evocare direttamente una
        creatura cavalcabile.\footnote{Il docente afferma che Nintendo avrebbe
        perso; la vicenda legale tra Nintendo, The Pokémon Company e
        Pocketpair era ancora in corso di cronaca al momento della lezione, e
        se ne riporta l'affermazione così com'è.} Lo si può aggirare in modo
        molto semplice: la creatura non entra nella sfera, entra nello zaino o
        nelle scarpe; e quando la si lancia non la si cavalca, ma ci si infila
        ``come un canguro'' nel suo marsupio. Non si può brevettare una
        meccanica estremamente specifica: ``nessuna agenzia ti darà l'ok''.
  \item \textbf{Il Nemesis System} di \emph{Middle-earth: Shadow of Mordor},
        un ``caso triste'' che il docente cita spesso, perché per chi l'ha
        giocato è probabilmente la cosa più memorabile del gioco, e non lo si
        usa più. L'orda di Mordor (orchi, sul modello del \emph{Signore degli
        Anelli}) ha una \textbf{gerarchia interna}: capi, gruppi, sottogruppi.
        Quando si uccide, o non si uccide, un orco, questo può salire o
        scendere di grado. Se lo si uccide viene rimosso e qualcuno prende il
        suo posto; se sopravvive a uno scontro, si ricorderà del giocatore e
        vorrà vendicarsi. Se uccide il giocatore, questi rinasce normalmente:
        la morte è prevista nell'esperienza come prosecuzione del sistema, e
        l'orco viene promosso per aver ucciso una figura di alto rango. Questo
        sistema è \textbf{brevettato}. Il docente pensa che si possa fare
        qualcosa che faccia percepire lo stesso responso emotivo, ma è
        difficile.
\end{itemize}
Sono casi rari: in generale è molto liberatorio poter prendere qualsiasi
meccanica e ``rivestirla''.

\paragraph{Perché il design dominante si riproduce.} Il discorso di
\emph{Dark Souls} si ripete un po' perché funziona estremamente bene a
livello di design (si capirà perché), un po' perché \textbf{il target audience
se lo aspetta}. Un giocatore di soulslike che prova un altro soulslike si
aspetta quel sistema: glielo si può dare tranquillamente, e anzi si fa meno
lavoro, perché lo si riprende uno a uno. Sempre che non si voglia innovare:
\textbf{innovare costa}, e non è nemmeno una necessità.

\subsection{Gli unique selling point}

Gli \textbf{unique selling point} (o \emph{unique selling proposition},
USP) vanno valutati \textbf{prima della produzione}. Sono cose uniche: una
\textbf{commistione di fattori} che solo quel gioco ha. Di solito sono
\textbf{tre, al massimo tre}. Per \emph{Palworld}, ad esempio, il fatto che
ci sia un monster collectible con un'arte ``pucciosa'' e le armi: non è solo
questo, ma è chiaramente l'elemento unico di punta.

\paragraph{Un nome famoso è un USP?} Valerio chiede se possa essere un USP
anche qualcosa che non è una feature del gioco: ad esempio, il nome di
Kojima. La risposta è ``sni''. Volendo essere molto generosi, sì; ma è un
\emph{selling point}, non uno \emph{unique} selling point, perché di giochi
``di Kojima'' ce n'è già più di uno. Certo, se si è Kojima si va da Sony o da
Microsoft (``mi paga uno, mi paga l'altro, è uguale'') e si ottengono i
finanziamenti. Ma non bisogna limitarsi a questo: il docente è convinto che,
nella direzione in cui sta andando l'industria, non sarà più così vero.

\paragraph{Il terzo capitolo di una saga.} Uno studente racconta che, al primo
anno, si era discusso di quale fosse l'USP di un \emph{Borderlands}: si
citavano lo shooter e molte altre cose, finché qualcuno disse ``è il terzo''.
Il punto è: che cosa dà in più rispetto agli altri? Il docente ricorda il
caso del \emph{Pre-Sequel}, che è insieme un prequel e un sequel: quello sì è
un modo unico di vendere il prodotto. Allo stesso modo, in \emph{Hollow
Knight: Silksong} l'USP non è ``siamo nel mondo di \emph{Hollow Knight}'', ma
``non gioco con il Cavaliere, gioco con Hornet''. Si può dire che il terzo
titolo di una saga è unico, ma bisogna \textbf{formularlo in maniera più
chiara}: il fatto che il titolo faccia parte della saga e \emph{aggiunga un
pezzo}. Lo si fa con gli spin-off: ambientato in quel mondo, ma con queste
meccaniche in più. Il nuovo \emph{A Plague Tale} citato in aula mantiene la
mitologia e la proprietà intellettuale (IP, \emph{intellectual property}) della
serie, ma ha meccaniche diverse: meno stealth, più action.

\paragraph{Una parentesi: fate domande.} A uno studente che si scusa per aver
riportato la lezione indietro, il docente risponde che è assolutamente
legittimo: ``non ho interesse a concludere la lezione, fare le tre ore e
andare a casa. Siamo qua per discutere''. Se c'è qualcosa che non si capisce,
ci si ferma e se ne parla; ``e sono pure pagato per parlare di giochi''. Se le
domande deviano il discorso e rendono le informazioni più utili, tanto meglio.

\paragraph{Esempio: gli USP di Enotria.} Essere un soulslike è un
\emph{selling point} (\emph{Enotria} vende anche perché è un soulslike), ma
non è \emph{unique}. Gli USP di \emph{Enotria} sono altri; il docente fatica a
trovarne un terzo, ma due li trova:
\begin{enumerate}
  \item sistemi da soulslike con una direzione artistica lontana dal
        grimdark: è il \textbf{Summer Souls}, e in più è \textbf{ambientato in
        Italia}, con il folklore italiano che si fonde con la mitologia
        tipica dei soulslike, boss enormi e mostruosi;
  \item il \textbf{sistema di maschere}: in un contesto soulslike, si possono
        avere diverse \emph{build} e cambiarle in tempo reale durante il
        combattimento, come se in \emph{Bloodborne} si cambiasse arma mentre
        si combatte (cosa che, più o meno, in parte si può fare).
\end{enumerate}
Si possono vendere tante cose, ma per pitchare il gioco come unico bisogna
far emergere questo. Per giochi piccoli l'USP potrebbe essere anche uno solo
(ma il docente invita a identificarne di più); per giochi molto grandi in
genere si arriva a tre. Sono cose che nessun altro ha fatto: in quella forma,
in quella dimensione, ``in quella salsa'', solo il proprio gioco.

\paragraph{Esempio: Metroid Other M e Metroid Dread.} Il docente, grande fan
della saga, porta l'esempio di \emph{Metroid}. Gli USP si recuperano anche da
trailer e interviste.
\begin{itemize}
  \item \emph{Metroid: Other M} è l'ennesimo capitolo della saga, e il suo USP
        era Samus come protagonista \textbf{umana}: si racconta il suo
        passato in un contesto in cui Samus parla e si relaziona con gli
        altri. Samus è riconoscibile, amata dai fan, ma più umana, e parla:
        ``infatti è piaciuta a pochissimi''. Ha delle intuizioni interessanti,
        ma la narrativa non è il suo punto forte. Samus, in genere, non parla
        in nessun gioco; qui ``parla e straparla'', mostrando lati deboli che
        non sono emersi nella maniera ideale per il target.
  \item \emph{Metroid Dread} lo ha capito: Samus parla, ma non in una lingua
        comprensibile: ma nella lingua degli alieni (i Chozo, da
        cui è stata cresciuta). È ``la cosa più figa dell'universo'': non
        mostra paura, è dura, cruda, precisa, esattamente la cacciatrice di
        taglie che si è abituati a giocare.
\end{itemize}
Ma l'USP di punta di \emph{Dread} è un altro: gli \textbf{E.M.M.I.}, nemici
che inseguono e cacciano il giocatore e che non si possono uccidere fino a un
certo punto. Ed è un collegamento con la storia della serie: in
\emph{Metroid Fusion} c'era già un nemico simile, l'SA-X, una copia della
tuta di Samus presa vita a causa di un parassita, che dà la caccia al
giocatore in alcune sezioni del livello. Lì era estremamente limitato (il GBA
era un'altra epoca); gli E.M.M.I. di \emph{Dread} sono quel concetto espanso,
e qui è molto più predominante. L'USP è simile, ma diventa il centro
dell'esperienza.

\subsection{Che cosa resta da vedere}

La prossima volta si concluderà questa parte e si vedrà che cos'altro si può
fare con i competitor, perché non è limitato a questo. Questi sono modi di
analizzarli ad alto livello, ponendosi un certo tipo di domande; ma molti dei
framework di design visti servono anche a \textbf{spacchettare i giochi}, e
conviene usarli nella fase di analisi.

% ---------------------------------------------------------------------
\section{Il task della settimana}
\label{sec:task-gd4}

Il nuovo task viene assegnato a metà lezione, subito dopo il production
dictionary (sezione~\ref{sec:production-dictionary}). La scadenza è
\textbf{lunedì della settimana successiva} (``lunedì tra una settimana'').

\begin{description}
  \item[Che cosa] Il \textbf{production dictionary} di un qualsiasi gioco
        \textbf{Pokémon della saga principale} (\emph{mainline}) a scelta. Non
        valgono gli spin-off: niente \emph{Mystery Dungeon}, niente
        \emph{Pokémon Snap}, niente \emph{Pokémon Rumble} (``ce ne sono un
        sacco, non li voglio''). A Daniele, uno studente, viene assegnato
        invece un \emph{Final Fantasy} con un numero inferiore a 10.
  \item[Dove] Questa volta spetta agli studenti ``l'onere e l'onore'' di
        creare la cartella e organizzare tutto secondo le regole del Drive
        (sezione~\ref{sec:drive-regole}).
  \item[Strumento] Un documento di Google, con le regole formali del corso
        (sezione~\ref{sec:regole-formali}). L'impostazione della pagina resta
        quella standard, salvo indicazioni diverse: sempre A4, stampabile.
  \item[Vincolo] \textbf{Non si deve superare una pagina}. Il limite serve a
        far ragionare sulla \textbf{priorità} delle parole e dei termini
        necessari: quali vanno spiegati per primi, perché più fraintendibili o
        più necessari alla comprensione.
\end{description}

Uno studente fa notare che i termini sono tanti e che, con limiti del genere,
si rischia di dover correre: proprio per questo il limite c'è. Se ne
riparlerà quando tutti avranno consegnato e il docente avrà letto i
documenti.

Valgono, come sempre, le golden rules (sezione~\ref{sec:golden-rules}) e i
termini banditi (sezione~\ref{sec:vocabolario}); e vale quanto detto sui
template (sezione~\ref{sec:due-principi}): nella lezione non viene mostrato un
esempio di production dictionary proprio perché non ce n'è uno da copiare.

\paragraph{Prossimi appuntamenti.}
\begin{itemize}
  \item Le slide della lezione non vengono caricate subito: lo saranno quando
        la lezione sarà conclusa, cioè \textbf{mercoledì}.
  \item \textbf{Venerdì} la lezione si terrà \textbf{da remoto}, a causa
        dello sciopero dei mezzi (lo si scriverà anche su Discord per chi è
        arrivato dopo).
  \item La dimostrazione di formattazione di un documento si farà in presenza,
        \textbf{venerdì 6} (sezione~\ref{sec:recap-documentazione}).
\end{itemize}

% ---------------------------------------------------------------------
\section*{Punti chiave del capitolo}
\addcontentsline{toc}{section}{Punti chiave del capitolo}

\begin{itemize}
  \item L'\textbf{atomo di gioco} di Koster: problema, preparazione,
        \textbf{core mechanic} calata in uno spazio, \textbf{feedback},
        successo o fallimento (con ricompensa), e infine l'aggiornamento del
        \textbf{modello mentale} del giocatore, che vive solo nella sua testa.
  \item La core mechanic è l'azione specifica che il giocatore ripete più
        spesso; il verbo di gioco, in senso gergale, è più generico.
  \item Il feedback \textbf{non è mai abbastanza}: se qualcosa non viene
        capito, spesso il problema è lì. Si itera: si osserva, si rinforza, si
        ritesta.
  \item Il \textbf{production dictionary} chiarisce i termini necessari a
        comprendere il gioco; di solito è un documento a sé, ha corrispettivi
        in altri dipartimenti (il \emph{lexicon}) e richiede
        \textbf{priorità}: prima i termini più fraintendibili o necessari.
  \item La documentazione non sostituisce il gioco, non tutti i documenti sono
        sempre necessari, e \textbf{non esistono template}: si prende il
        meglio, lo si ibrida, lo si adatta.
  \item Il design è il \textbf{miglior compromesso possibile} tra team e
        know-how, visione e scienza, mercato e target audience, produzione e
        soldi; è \textbf{50\% scienza e 50\% arte}.
  \item ``Senza dati si è solo un'altra persona con un'opinione'': servono
        dati, quantitativi e qualitativi, da almeno un centinaio di persone,
        raccolti in pre-produzione, osservati, filmati, ordinati in archetipi e
        correlati.
  \item Il design thinking in cinque passi: imparare dalle persone, trovare
        schemi, definire principi, rendere tangibile, iterare senza sosta.
  \item La mappa della ricerca: client pitch direction analysis,
        stakeholder, problema contro necessità, competitor, utenti, mercato,
        gameplay, fino a una \textbf{visione} che non cambia in corsa.
  \item I \textbf{competitor diretti} condividono la maggior parte dei punti
        insindacabili; gli \textbf{indiretti} uno solo, molto attinente; è uno
        spettro. I giochi in \textbf{opposizione} aiutano a innovare.
  \item Il mercato può essere un \textbf{red ocean} o un \textbf{blue ocean}, e
        ha regole (Cina, Germania, Belgio) a cui pensare prima.
  \item Strumenti di indagine: cinque W, nove perché e root cause analysis,
        interviste strutturate, must have contro nice to have, scenari
        estremi.
  \item Gli stakeholder si classificano per \textbf{potere e interesse}
        (promotori, latenti, difensori, apatici) o per \textbf{potere,
        legittimità e urgenza}.
  \item Il \textbf{problema} sta in una struttura esistente e si risolve in
        fretta (un cerotto); la \textbf{necessità} richiede una struttura nuova
        (i punti di sutura). Ishikawa, Liberating Structures e Cynefin aiutano
        a categorizzare.
  \item Nell'analisi dei competitor: SWOT, design dominante (le meccaniche,
        salvo rari brevetti, sono di tutti), \textbf{unique selling point}
        (al massimo tre, davvero unici).
  \item Task: production dictionary di un Pokémon mainline, \textbf{una sola
        pagina}, per lunedì della settimana successiva.
\end{itemize}
